\documentclass[10pt,a4paper]{amsart}
\usepackage{amsmath,amssymb,amsthm,graphicx}
\usepackage{hyperref}
\numberwithin{equation}{section}

\usepackage[mathscr]{euscript}
\usepackage[all]{xy}
\usepackage{ytableau}
\usepackage{enumitem}

\DeclareMathOperator{\ch}{ch}
\DeclareMathOperator{\gr}{gr}
	
\newcommand{\A}{{\mathbb{A}}}
\newcommand{\Ad}{\mathrm{Ad}}
\newcommand{\ad}{\mathrm{ad}}
\newcommand{\Addable}{\mathcal{A}}
\newcommand{\af}{{\mathrm{af}}}
\newcommand{\Ainf}{{A_\infty}}
\newcommand{\al}{\alpha}
\newcommand{\alv}{\alpha^\vee}
\newcommand{\alt}{\mathrm{alt}}
\newcommand{\Aut}{\mathrm{Aut}}
\newcommand{\Ax}{{\tilde{\mathbb{A}}}}
\renewcommand{\b}{\bar{b}}
\newcommand{\bl}{\mathbf{l}}
\newcommand{\BFM}{\mathrm{BFM}}
\newcommand{\bGr}{\overline{\Gr}}
\newcommand{\ba}{\bar{b}}
\newcommand{\bFl}{\overline{\Fl}}
\newcommand{\Bv}{B^\vee} % Borel of \Gv
\newcommand{\C}{\mathbb{C}}
\newcommand{\Cinf}{{C_\infty}}
\newcommand{\cl}{\mathrm{cl}}
\newcommand{\co}{\xi}
\newcommand{\col}[2]{{\color{#1}{#2}}}
\newcommand{\cT}{\mathcal{T}}
\newcommand{\cZ}{\mathscr{Z}}
\newcommand{\cZGv}{\cZ_{G^\vee}}
\newcommand{\eps}{\varepsilon}
\newcommand{\fb}{\mathfrak{b}}
\newcommand{\fbv}{\mathfrak{b}^\vee}
\newcommand{\fg}{\mathfrak{g}}
\newcommand{\fgv}{\mathfrak{g}^\vee}
\newcommand{\fgvs}{\mathfrak{g}^{\vee*}}
\newcommand{\fh}{\mathfrak{h}}
\newcommand{\fhv}{\fh^\vee}
\newcommand{\Fl}{\mathcal{F}l}
\newcommand{\fn}{\mathfrak{n}}
\newcommand{\fnv}{\mathfrak{n}^\vee}
\newcommand{\Fac}{\delta}
\newcommand{\Frac}{\mathrm{Frac}}
\newcommand{\frg}{\fg}
\newcommand{\fso}{\mathfrak{so}}
\newcommand{\fsp}{\mathfrak{sp}}
\newcommand{\Gab}{\Gamma_\bullet}
\newcommand{\GL}{\mathrm{GL}}
\newcommand{\Gabh}{\hat{\Gamma}_\bullet}
\newcommand{\Gad}{G_\mathrm{ad}}
\newcommand{\GKM}{\Psi}
\newcommand{\Gr}{\mathrm{Gr}}
\newcommand{\Gro}[1]{\mathfrak{G}^Q_{#1}}
\newcommand{\Gs}{\Gamma^*} % temporary
\newcommand{\Gv}{G^\vee} % Langlands dual group of G
\newcommand{\he}{\hat{e}}
\newcommand{\hG}{\hat{\Gamma}}
\newcommand{\hGKM}{\hat{\GKM}}
\newcommand{\hh}{\hat{h}}
\newcommand{\hq}{\hat{q}}
\newcommand{\hw}{\hat{w}}
\newcommand{\hLa}{\hat{\Lambda}}
\newcommand{\Hom}{\mathrm{Hom}}
\newcommand{\hs}{\hat{s}}
\newcommand{\hP}{\hat{P}}
\newcommand{\hQ}{\hq}
\newcommand{\hS}{{\hat{S}}}
\newcommand{\hz}{\hat{z}}
\newcommand{\Iaf}{\tilde{I}}
\newcommand{\id}{\mathrm{id}}
\newcommand{\II}{I_Z}
\newcommand{\jpi}{\Xi} % the special element corresponding to j_\pi
\newcommand{\Kloc}{\mathbb{K}_{\mathrm{loc}}}
\newcommand{\K}{\mathbb{K}}
\newcommand{\Kx}{\tilde{\mathbb{K}}}
\newcommand{\la}{\lambda}
\newcommand{\La}{\Lambda}
\newcommand{\LGr}{\mathrm{LGr}}
\newcommand{\Lie}{\mathrm{Lie}}
\newcommand{\loc}{\mathrm{loc}}
\newcommand{\LL}{\mathbb{L}}
\newcommand{\Lx}{\tilde{\mathbb{L}}}
\newcommand{\m}[1]{\overline{#1}}
\newcommand{\Map}{\mathrm{Map}}
\newcommand{\nilA}{\A_\af} % affine nil Hecke algebra
\newcommand{\nilAx}{\Ax_\af}
\newcommand{\om}{\varpi}
\newcommand{\omv}{{\varpi^\vee}}
\newcommand{\pair}[2]{\langle #1 , #2\rangle}
\newcommand{\Pet}{Z_{\nilA}(S)} % Peterson algebra
\newcommand{\Petx}{\tilde{Z}_{\nilA}(S)}
\renewcommand{\O}{\mathscr{O}}
\newcommand{\pt}{\mathrm{pt}}
\newcommand{\Px}{\mathbb{L}_e}
\newcommand{\PGL}{\mathrm{PGL}}
\newcommand{\phit}{\tilde{\phi}}
\newcommand{\Q}{\mathbb{Q}}
\newcommand{\QRep}{Q(T)}
\newcommand{\Quo}{A_n}
\newcommand{\Quom}{A_n^{\mathfrak{m}}}
\newcommand{\rat}{\mathrm{rat}}
\newcommand{\reg}{\Delta}
\newcommand{\R}{R_n} % ring generated by single rows
\newcommand{\Rep}{{R(T)}}
\newcommand{\Removable}{\mathcal{R}}
\newcommand{\Schub}{\mathfrak{S}}
\newcommand{\SL}{\mathrm{SL}}
\newcommand{\SO}{\mathrm{SO}}
\newcommand{\Spec}{\mathscr{S}}
\newcommand{\ks}{g}
\newcommand{\kst}{\tilde{g}}
\newcommand{\kt}[2]{\tilde{g}^{(#1)}_{#2}}
\newcommand{\SP}{\mathscr{SP}}
\newcommand{\Sp}{\mathrm{Sp}}
\newcommand{\Sym}{\mathrm{Sym}}
\newcommand{\sym}{\mathrm{sym}}

\renewcommand{\t}{\xi}
\newcommand{\resC}[1]{C_{#1}}
\newcommand{\trace}{\mathrm{tr}}
\newcommand{\trsp}[1]{{}^t\!{#1}} % transpose
\newcommand{\Trsp}[1]{{}^T\!{#1}}
\newcommand{\tX}{\tilde{X}}
\newcommand{\ul}[1]{\underline{#1}}
\newcommand{\uf}{u}
\newcommand{\Wafgr}{\hat{W}_G^0} % affine Grassmannian elements
\newcommand{\wv}{\varpi^\vee}
\newcommand{\Wxgr}{{\Wx}^0} % affine Grassmannian elements
\newcommand{\Wx}{\tilde{W}_\af}
\newcommand{\Wxx}{\hat{W}_\af}
\newcommand{\wz}{w_\circ}
\newcommand{\Z}{\mathbb{Z}}
\newcommand{\de}{\varrho}
\newcommand{\sinf}{\frac{\infty}{2}}
\newcommand{\fac}{\psi}

\newcommand{\ha}{f}
\usepackage{chngcntr}

\def\NN{\mathord{\mathbb{N}}}
\def\CC{\mathord{\mathbb{C}}}
\def\QQ{\mathord{\mathbb{Q}}}
\def\RR{\mathord{\mathbb{R}}}
\def\ZZ{\mathord{\mathbb{Z}}}
\def\PP{\mathord{\mathbb{P}}}
\def\LL{\mathbb{L}}
\def\Tableau#1{\ygtab{#1}}
\newcommand{\bm}[1]{\mbox{\boldmath{$#1$}}}

\newtheorem{Theorem}{Theorem}[section]
\newtheorem*{Theorem*}{Theorem}
\newtheorem{Corollary}[Theorem]{Corollary}
\newtheorem{Lemma}[Theorem]{Lemma}
\newtheorem{Proposition}[Theorem]{Proposition}

\theoremstyle{definition}
\newtheorem{Definition}[Theorem]{Definition}
\newtheorem{Note}[Theorem]{Note}
\newtheorem{Example}[Theorem]{Example}
\newtheorem{Remark}[Theorem]{Remark}

\title{The relativistic Toda substitution identifies quantum and affine equivariant K-theory Schubert bases}
\author{Takeshi Ikeda, Shinsuke Iwao, Satoshi Naito, Kohei Yamaguchi}
\date{doi:10.3842/SIGMA.2026.065; CC BY 4.0}
\begin{document}
\maketitle
\noindent\textit{Allowed standard prerequisites (editorial).}\begin{itemize}
\item Equivariant K-theory, Demazure/DAHA operators and standard symmetric-function theory. Original complete bibliography and source-local definitions.
\end{itemize}\medskip
\subsection[K-theoretic Peterson map: abstract form]{$\boldsymbol{K}$-theoretic Peterson map: Abstract form}
To describe the $K$-Peterson map more precisely, we fix some notation of the affine Weyl group.
Let $W_G$ be the Weyl group of $(G,T)$, and $\hat{W}_G=W\ltimes Q^\vee$ the affine Weyl group,
 where $Q^\vee$ is the coroot lattice.
 Let $\hat{W}_G^0$ be the set of minimal-length coset representatives for
 $\hat{W}_G/W_G$. For~\smash{$x\in \hat{W}_G^0$}, there is
 an associated Schubert structure sheaf $\mathscr{O}_x$ in \smash{$K_*^T(\Gr_{G})$}.
 These sheaves form an $\Rep$-basis of $K_*^T(\Gr_{G})$, where $\Rep$ denotes the representation ring of $T$.
On the quantum side, we have the Schubert class $\mathscr{O}^w$ for each $w\in W_G$.
 For \smash{$x\in \hat{W}_G^0$},
 write it as~${x=wt_\xi}$ with~${w\in W}$, $\xi\in Q^\vee$,
 where $t_\xi$ is the translation element corresponding to $\xi\in Q^\vee$.
 Let $QK_T(G/B)_Q$ denote the localization of $QK_T(G/B)$ by
 the Novikov variables $Q_1,\dots,Q_r$, where~$r$ is the rank of $G$.
 The $K$-Peterson map, at the abstract level of Schubert bases,
maps $Q^\xi\mathscr{O}^w \in QK_T(G/B)_Q$ to $\mathscr{O}_{x}\in K_*^T(\Gr_G)$, where $Q^\xi$ denote the product of the Novikov variables corresponding to $\xi$. This map establishes a~correspondence between the quantum and affine Schubert calculus.

\subsection[K-Peterson map from the Relativistic Toda lattice]{$\boldsymbol{K}$-Peterson map from the Relativistic Toda lattice}

Let us explain the key idea of our construction before going into the details.
By solving the relativistic Toda lattice, we obtain a~birational map between the phase space $\mathscr{Y}$ of the relativistic Toda lattice
and a~certain centralizer family $\mathscr{Z}$ associated with $\PGL_n(\C)$.
Both varieties are defined over $T$, the maximal torus of $\SL_n(\C)$, and their coordinate rings have the following geometric interpretations:
\begin{equation}\label{eq:O_geom}
\mathscr{O}(\mathscr{Y})
\cong
QK_T^{\mathrm{pol}}(\SL_n(\C)/B),\qquad
\mathscr{O}(\mathscr{Z})
\cong
K_*^T(\Gr_{\SL_n}).
\end{equation}
Here $QK_T^{\mathrm{pol}}(\SL_n(\C)/B)$ is the polynomial version of the
quantum $K$-ring of the flag manifold (see Section~\ref{sec:QK}).
The second isomorphism of~\eqref{eq:O_geom} was proved in~\cite{ISY}, which is a~$K$-theory analogue of a~result for $H_*^T(\Gr_G)$ due to Peterson~\cite{Pet} and Ginzburg~\cite{Gi} independently.
We obtain the following commutative diagram (see Theorem~\ref{thm:main})
\begin{equation}
\label{eq:comdiag}
\begin{split}&
\xymatrix{
QK_T^{\mathrm{pol}}(\SL_n(\C)/B)_{Q}
\ar[d] \ar[r] \ar@{}[dr] & K_*^T(\Gr_{\SL_n})\big[{\sigma_i^{-1},\tau_i^{-1}\mid 1\leq i\leq n}\big]
 \ar[d] \\
\mathscr{O}(\mathscr{Y}^\circ) \ar[r]_{\widetilde{\Phi}_n} & {\mathscr{O}}(\mathscr{Z}^\circ),
}\end{split}\end{equation}
 where the top arrow is Kato's map, and the bottom arrow is defined by~\eqref{eq:tPhi} below.
In this diagram, all maps are isomorphisms, and $\mathscr{Y}^\circ$ and $\mathscr{Z}^\circ$ are certain open dense subsets of $\mathscr{Y}$ and~$\mathscr{Z}$ respectively, and $\sigma_i$, $\tau_i$ are discussed in detail in Section~\ref{sec:tau}.
This perspective extends the work~\cite{LS:double-Kostant,LS10} by Lam and Shimozono for
(co)homology and the classical Toda lattice.

\subsection[K-theoretic double k-Schur function]{$\boldsymbol{K}$-theoretic double $\boldsymbol{k}$-Schur function}
\label{sec:double_K-k}

Another ingredient of
our work is the \emph{$K$-theoretic double $k$-Schur function} introduced in~\cite{ISY}, which is an equivariant deformation of the $K$-theoretic $k$-Schur function introduced by Lam, Schilling, and Shimozono~\cite{LSS:K}.
The representation ring $\Rep$ is given as
\[
\Z\big[{\rm e}^{\pm a_i} (1\le i\le n)\big]/\bigl({\rm e}^{a_1+\cdots+a_n}-1\bigr).
\]
For \smash{$x\in \hat{W}_G^0$}, the $K$-theoretic double $k$-Schur function \smash{$g_x^{(k)}(y|b)$} is a~symmetric formal power series in the infinitely many variables $y=(y_1,y_2,\dots)$ with coefficients in~$\Rep$, where we set $k=n-1$.
It depends on the sequence $b=(b_1,\dots,b_n)$ of equivariant parameters, where each~$b_i$~is identified
with $1-{\rm e}^{-a_i}\in \Rep$.
 Let \smash{$\hat{\Lambda}^{\Rep}_{(n)}$} be the $\Rep$-span of \smash{$g_x^{(k)}(y|b)$}, $
 x\in \hat{W}_G^0$. Then we have an isomorphism
\[
 K_*^T(\Gr_{\SL_n})\cong \hat{\Lambda}^{\Rep}_{(n)}
\]
 of $\Rep$-algebras such that the structure sheaf~$\mathscr{O}_x$ corresponds to the \emph{closed} $K$-theoretic double $k$-Schur function~\smash{$\tilde{g}_x^{(k)}(y|b):=\sum_{z\le x}g_z^{(k)}(y|b)$}, where $\le $ denotes
 the Bruhat order on $\hat{W}_G^0$.

There is a~bijection $\hat{W}_G^0\cong \mathscr{P}^{(k)}$, where
$\mathscr{P}^{(k)}$ denotes the set of $k$-bounded partitions, i.e., the~partition $\la=(\la_1,\dots,\la_i)$ such that
$\la_1\le k$.
If \smash{$x\in \hat{W}_G^0$} corresponds to
\smash{$\la\in \mathscr{P}^{(k)}$},
then we write \smash{$g_x^{(k)}(y|b)$} \big(resp.\ \smash{$\tilde{g}_x^{(k)}(y|b)$}\big)
as \smash{$g_\la^{(k)}(y|b)$} \big(resp.\ \smash{$\tilde{g}_\la^{(k)}(y|b)$}\big).

We have derived determinantal formulas for
\smash{$g^{(k)}_\la(y|b)$} and \smash{$\tilde{g}_\la^{(k)}(y|b)$} (see Theorems~\ref{thm:determinant_formula_for_ksmall_g} and~\ref{thm:determinantal_formula_for_g})
for a~$k$-\emph{small} $k$-bounded partition $\la$; a~partition $\la\in \mathscr{P}^{(k)}$ is said to be $k$-small if $\la_1+\ell(\la)\le n$, where
$\ell(\la)$ is the number of nonzero parts of $\la$.

\subsection[k-rectangles and the tau-functions]{$\boldsymbol{k}$-rectangles and the $\boldsymbol{\tau}$-functions}
The so-called $\tau$-\emph{functions} in the theory of integrable systems also play an important role in our geometric context. For $1\le i\le n$, we define
$\tau_i,\sigma_i\in \mathscr{O}(\mathscr{Z})$ as the $i$-th principal minor determinants of certain matrices
related to the centralizer family (see Section~\ref{sec:tau}).

For $1\le i\le k$, let $R_i$ denote the partition $(i)^{n-i}\in \mathscr{P}^{(k)}$ of rectangular shape, and $R_n=\varnothing$.
Note that each $R_i$ $(1\le i\le n-1)$ is a~maximal $k$-small $k$-bounded partition.
Under the isomorphism
\smash{$ \mathscr{O}(\mathscr{Z})
\cong \hat{\Lambda}^{\Rep}_{(n)}$},
$\tau_i$ and $\sigma_i$
correspond to \smash{$g_{R_{n-i}}^{(k)}(y|b)$} and \smash{$\tilde{g}_{R_{n-i}}^{(k)}(y|b)$} up to
some simple factors respectively (see Corollary~\ref{cor:k rect and tau}).

These functions are fundamental because
\smash{$\tilde{g}_{\la}^{(k)}(y|b)$} satisfies the $k$-rectangle factorization property (see Theorem~\ref{thm:$k$-rect})
\[
\tilde{g}_{R_i\cup \la}^{(k)}(y|b)=\tilde{g}_{R_i}^{(k)}(y|b)
\tilde{g}_{\la}^{(k)}\bigl(y|\omega^ib\bigr),
\]
where $\omega$ is the permutation sending $b_i$ to $b_{i+1}$ with $b_{n+1}=b_1$.
Thanks to this result, we only need to study the functions \smash{$\tilde{g}_\la^{(k)}(y|b)$} associated with the
$k$-\emph{irreducible} $k$-bounded partition $\la$; i.e., an element of $\mathscr{P}^{(k)}$ which is not expressed as $R_i\cup\mu$ for $1\le i\le n-1$ and $\mu\in \mathscr{P}^{(k)}$
with~${\mu\ne \varnothing}$.

\subsection{Quantum double Grothendieck polynomials}

For $w\in S_n$, let $\Gro{w}(z| \eta)$ be the \emph{quantum double
Grothendieck polynomial} due to Lenart and Maeno~\cite{LM}.
This is a~polynomial in two sets of variables
$z_1,\dots,z_n$ and $\eta_1,\dots,\eta_n$ with coefficients in
$\Z[Q_1,\dots,Q_{n-1}]$.
We basically follow the notation in~\cite{MNS2}, however,
there are some differences in the identification of the equivariant parameters.
In particular, $\eta_i$ is identified with $1-{\rm e}^{a_{n-i+1}}\in \Rep$ (see~\eqref{eq:eta-b} for more details).
We denote $\Rep[z_1,\dots,z_n,Q_1,\dots,Q_{n-1}]$ by $\Rep[z,Q]$.

In the context of the relativistic Toda lattice,
$z_i$, $Q_i$ are interpreted as dynamical variables (see Section~\ref{sec:Rel Toda}).
There are conserved quantities $F_i(z,Q)\in \Rep[z,Q]$ $(1\le i\le n)$ of the relativistic Toda lattice.
Let $\mathscr{I}_n^{Q,\mathrm{pol}}$ be the
ideal generated by $F_i(z,Q)-e_i({\rm e}^{-a_1},\dots,{\rm e}^{-a_n})$ $(1\le i\le n)$, where $e_i$ denotes the $i$-th elementary
symmetric polynomial.
The ring \smash{$\Rep[z,Q]/\mathscr{I}_n^{Q,\mathrm{pol}}$} is by our definition
the ring of regular functions $\mathscr{O}(\mathscr{Y})$ on the
phase space $\mathscr{Y}$.

Let $\mathscr{Y}^\circ$ be the open set of $\mathscr{Y}$ defined as the complement of the divisor given by the equation $Q_1\cdots Q_{n-1}=0$.
Due to~\cite{MNS2}, the ring $\mathscr{O}(\mathscr{Y}^\circ)=\mathscr{O}(\mathscr{Y})\big[Q_i^{-1}\mid 1\leq i\leq n-1\big]$
is identified with~\smash{$QK_T^{\mathrm{pol}}(\SL_n(\C)/B)_Q$}
and \smash{$\mathfrak{G}_w^Q(z| \eta)$} represents the Schubert structure sheaf \smash{$\mathscr{O}^w$} (see Remark~\ref{rem:QK=O(Y)} for more details).

\subsection{Correspondence of Schubert bases}
Following an analogous approach to that used in Kostant's construction of solutions to the Toda lattice, we obtain a~map
\[\Phi_n\colon \ \mathscr{O}(\mathscr{Y}^\circ)
\longrightarrow \hat{\Lambda}^{\Rep}_{(n)}\big[\tau_i^{-1},\sigma_i^{-1}\mid 1\leq i\leq n\big]
\cong\mathscr{O}(\mathscr{Z}^\circ)\]
of $\Rep$-algebras as
\begin{equation}
z_i\mapsto \frac{\tau_i\sigma_{i-1}}{\sigma_i \tau_{i-1}},\qquad
Q_i\mapsto \frac{\tau_{i-1}\tau_{i+1}}{\tau_i^2}.\label{eq:Phi}
\end{equation}

Let $\omega_k$, with $k=n-1$, be an involution on $\hat{W}_G^0$ called the $k$-\emph{conjugation} defined by replacing~$s_i$ with $s_{n-i}$ for $i \in I:=\{1,\dots,n-1\}$ in any reduced expression of \smash{$x\in \hat{W}_G^0$}.
There is an automorphism $\sigma$ of \smash{$\hat{\Lambda}^{\Rep}_{(n)}$} as an~$\Rep$-algebra
sending $h_i(y)\ (i\in\mathbb{N})$ to $1+h_1(y)+\cdots+h_i(y)$, where~$h_i(y)$ is the
$i$-th complete symmetric function.
In particular, we have $\sigma(\tau_i)=\sigma_i$.
Define the map $\tilde{\Phi}_n$ as follows
\begin{equation}
\label{eq:tPhi}
\tilde{\Phi}_n=\sigma \circ \Phi_n.
\end{equation}
The main result of this paper is the following.

\begin{Theorem}\label{thm:main}
We have the commutative diagram~\eqref{eq:comdiag}. More precisely,
for \smash{$x\in \hat{W}_G^0$}, write $x=wt_\xi$ with $w\in W$, $\xi\in Q^\vee$. Then
\[
\tilde{\Phi}_n\bigl(
Q^\xi \Gro{w}(z| \eta)\bigr)=
\tilde{g}_{x^{\omega_k}}^{(k)}(y|b).
\]
\end{Theorem}

Thus, the substitution map ${\Phi}_n$ agrees with Kato's map up to a~twist, after identifying
 $QK_T(\SL_n(\C)/B)$ with $\mathscr{O}(\mathscr{Y})$ \big(resp.\ $K_*^T(\Gr_{\SL_n})$ with $\O(\mathscr{Z})$\big).

Let us outline the proof.
We begin by proving the theorem for $x=s_0$ in a~purely combinatorial manner (see Proposition~\ref{prop:key}).
In this proof, we utilize basic properties of the quantum double Grothendieck polynomials, the realization of $K_*^T(\Gr_{SL_n})$ as $\mathscr{O}(\mathscr{Z})$, and the explicit description of $\Phi_n$.
This step forms the technical core of the paper.
An essential idea to prove the general case is the use of the action of Demazure operators $D_x^Q$, $x\in \hat{W}_G^0$, on $QK_T(\SL_n(\C)/B)_Q$.

As a~straightforward consequence of the theorem for the specific case
$x=s_0$,
we demonstrate that the map $\tilde{\Phi}_n$
intertwines the actions of the Demazure operators on both the quantum and affine sides (see Corollary~\ref{cor:Phi_commutes_w_Dem}).
The final result needed to complete the proof is the
precise formula for the action of
$D_i^Q$ on the Schubert structure sheaves $\mathscr{O}^w$ in
$QK_T(\SL_n(\C)/B)$
(see Proposition~\ref{prop:left_Dem}).

\subsection[Factorization formula for the maximal k-irreducible partition nu\_n]{Factorization formula for the maximal $\boldsymbol{k}$-irreducible partition $\boldsymbol{\nu_n}$}\label{ssec:fac max}
As an application of
Theorem~\ref{thm:main}, we derive a~factorization formula of \smash{$\tilde
{g}_{\nu_n}^{(k)}(y|b)$}, where
\[
\nu_n=\bigcup_{i=1}^{n-2}(n-i-1)^i.
\]
This partition $\nu_n$ is important because it is the unique maximal $k$-irreducible $k$-bounded partition.
The non-equivariant version of the following result was conjectured in~\cite{IIM}, and it was proved by Blasiak, Morse, and Seelinger~\cite{BMS}.

Let $\lfloor x\rfloor$ denote the greatest integer less than or equal to $x\in \mathbb{R}$.
\begin{Theorem}\label{thm:max factor}
Let $n$ be even and write $n = 2m$. Then
\begin{align}
\label{eq:even}
\tilde{g}^{(k)}_{\nu_n}(y|b)
= \prod_{i=0}^{\lfloor (m-1)/2 \rfloor}
 \frac{\Omega(b_{m-1-2i} | y)}{\Omega(b_{m+2-2i} | y)}
 \prod_{i=1}^{n-2}
 \tilde{g}^{(k)}_{(n-i-1)^i}\bigl(y | \omega^{m+2i+1} b\bigr),
\end{align}
where
$\Omega(b_i | y)$ is defined in~\eqref{eq:s0_Omega}.
Let $n$ be odd. Then
\begin{align}
\label{eq:odd}
\tilde{g}^{(k)}_{\nu_n}(y|b)
= \prod_{i=1}^{n-2}
 \tilde{g}^{(k)}_{(n-i-1)^i}\bigl(y | \omega^{2i+1} b\bigr).
\end{align}
\end{Theorem}


Note that all the factors \smash{$\tilde{g}^{(k)}_{(n-i-1)^i}\bigl(y|\omega^{2i+1}b\bigr)$}
are expressed as determinants because $(n-i-1)^i$ is $k$-small.


\section{Relativistic Toda lattice and the centralizer family}\label{sec:RTL}
The aim of this section is to explain how the map $\Phi_n$ is
obtained by solving the relativistic Toda lattice.
Taking into account
recent developments,
we will review the construction in~\cite{IIM}.



\subsection{Relativistic Toda lattice}\label{sec:Rel Toda}
The \emph{relativistic Toda lattice equation}
was introduced by Ruijsenaars~\cite{Rui}.
In this paper, we start from the Lax equation due to Suris~\cite{Suris}.
Let
\begin{align}\label{eq:def_of_M_and_Q}
M=\begin{pmatrix}
z_1 & -1 & 0 & \cdots &0\\
 & z_2 & -1 & \ddots &\vdots\\
 & & \ddots & \ddots & 0\\
 & & &\ddots & -1\\
 & & & & z_n
\end{pmatrix},\quad
N=\begin{pmatrix}
1 & 0 & 0 & \cdots &0\\
-Q_1z_1 & 1 & 0 & \ddots &\vdots\\
 & -Q_2z_2 & \ddots & \ddots & 0\\
 & & \ddots&1& 0\\
 & & &-Q_{n-1}z_{n-1} & 1
\end{pmatrix}
\end{align}
and $L:=MN^{-1}$, which we call the Lax matrix.
We consider the system of partial differential equations
\begin{equation}\label{eq:rel_Toda}
{\partial L}/{\partial t_i}=\big[L,\bigl(L^i\bigr)_{<}\big]\qquad\text{for}\quad 1\le i\le n-1,
\end{equation}
where $\bigl(L^i\bigr)_{<}$ is the
strictly lower triangular part of $L^i$.
The equation is a~group version of the finite open Toda lattice.



\subsection[Quantum K-ring of the flag variety]{Quantum $\boldsymbol{K}$-ring of the flag variety}\label{sec:QK}
The integrable system is related to the quantum
$K$-ring of the flag variety $\SL_n(\C)/B$, analogously to a~result of
Givental and Kim~\cite{GK} for the quantum cohomology ring.

The conserved quantities $F_i(z,Q)$
of the system~\eqref{eq:rel_Toda} are given by
\[
\det(\zeta E-L)=\sum_{i=0}^n(-1)^i F_i(z,Q)\zeta^{n-i}.
\]
Explicitly, we have (see Appendix~\ref{sec:F})
\begin{equation}\label{eq:polynomialF}
F_i(z,Q)=
\sum_{\substack{J\subset \{1,\dots,n\}\\ |J|=i}}
\prod_{{j\in J,\, j+1\notin J}}(1-Q_j)\prod_{j\in J}z_j.
\end{equation}




Let $R(T)[\![Q]\!]$ denote the
ring of formal power series in the variables $Q_1,\dots,Q_{n-1}$ with coefficients in $R(T)$.
The quantum $K$-ring $QK_T(\SL_n(\C)/B)$ is a~commutative $R(T)[\![Q]\!]$-algebra.
\begin{Theorem}[Maeno, Naito, Sagaki~\cite{MNS1}]\label{thm:MNS1}
There exists an isomorphism of $R(T)[\![Q]\!]$-algebras
\[
QK_T(\SL_n(\C)/B)\cong
\Rep[\![Q]\!][z_1,\dots,z_n]/\mathscr{I}_n^Q,
\]
where the ideal $\mathscr{I}_n^Q$ is generated by the elements
\begin{equation}\label{eq:conserved}
F_i(z,Q)-e_i({\rm e}^{-a_1},\dots,{\rm e}^{-a_n})\qquad \text{for}\quad 1\le i\le n. \end{equation}
\end{Theorem}
\begin{Remark}
Our $a_i$ is $-\epsilon_i$ in~\cite{MNS1}.
\end{Remark}In this geometric
context, $z_i$ is related to the class of
the universal line bundle (see~\cite{MNS1} for details),
and $Q_i$ $(1\le i\le n-1)$ are the Novikov variables.

\begin{Remark}\label{rem:KPSZ}
It follows from \cite[Corollary~2]{GL} that the elements in~\eqref{eq:conserved} are contained in the defining ideal of $QK_T(\SL_n(\C)/B)$ (see also~\cite{ACT, MNS1} for details).
We also note that Koroteev--Puskar--Smirnov--Zeitlin~\cite{KPSZ} proved that the elements in~\eqref{eq:conserved} generate the defining ideal of $
QK_T^{\mathrm{QM}}\bigl(\SL_n(\C)/B\bigr)$,
which is the limit of quasimap quantum $K$-ring of $T^*(\SL_n(\C)/B)$ at $\hbar = \infty$.
Moreover, Huq--Kuruvilla~\cite{HK} established that this limit is isomorphic to $QK_T(\SL_n(\C)/B)$ (more generally, for partial flag varieties).
\end{Remark}

Let $\mathscr{Y}$ be the affine subscheme
of $\mathbb{A}^{2n-1}$ with coordinates
$Q_1,\dots,Q_{n-1}$, $z_1,\dots,z_n$
whose defining ideal is generated by the polynomials~\eqref{eq:conserved}.
$\mathscr{Y}$ is a~scheme over $T$ whose fibers are the isolevel sets.
The coordinate ring, i.e., the ring of regular functions $\mathscr{O}(\mathscr{Y})$ is
considered to be a~polynomial version \smash{$QK_T^{\mathrm{pol}}(\SL_n(\C)/B)$} of $QK_T(\SL_n(\C)/B)$.

\begin{Remark}
\label{rem:QK=O(Y)}
Let us denote the ideal of $\Rep[z,Q]$
generated by the polynomials~\eqref{eq:conserved}
by $\mathscr{I}_n^{Q,\mathrm{pol}}$.
The quotient ring
\smash{$\mathscr{O}(\mathscr{Y})=\Rep[z,Q]/\mathscr{I}_n^{Q,\mathrm{pol}}$} is
\smash{$QK_T^{\mathrm{pol}}(\SL_n(\C)/B)$} in the notation used in the introduction.
Although there are subtleties in relating \smash{$QK_T^{\mathrm{pol}}(\SL_n(\C)/B)$}
to $QK_T(\SL_n(\C)/B)$,
\smash{$QK_T^{\mathrm{pol}}(\SL_n(\C)/B)_Q=\mathscr{O}(\mathscr{Y}^\circ)$}
is actually a~subring of $QK_T(\SL_n(\C)/B)_Q$. Furthermore, the residue class of \smash{$\mathfrak{G}_w^Q$} in \smash{$QK_T^{\mathrm{pol}}(\SL_n(\C)/B)_Q$} corresponds to
$\mathscr{O}^w$ in $QK_{T}(\SL_n(\C)/B)$. For these facts, see \cite[Sections~4.1 and 4.2]{IIN}; it is straightforward to have the equivariant version based on \cite[Remark~6.3]{MNS1}.
\end{Remark}

If we consider
the special fiber $\mathscr{Y}_{\mathrm{uni}}$
corresponding to the case when the Lax matrix is \emph{unipotent}, and the corresponding
centralizer $\mathscr{Z}_{\mathrm{uni}}$, we obtain the non-equivariant $K$-Peterson isomorphism
$QK(\SL_n(\C)/B)_\loc^{\mathrm{pol}}\cong K_*(\Gr_{\SL_n})_\loc$ studied in~\cite{IIM}.
In fact, the element of $\mathscr{Z}_{\mathrm{uni}}$ is of the form
\[
\begin{pmatrix}
1 & h_1 & h_2 & \cdots &h_{n-1}\\
 & 1 & h_1 & \ddots &\vdots\\
 & & \ddots & \ddots & h_2\\
 & & &\ddots & h_1\\
 & & & & 1
\end{pmatrix},
\]
so $\mathscr{O}(\mathscr{Z}_\mathrm{uni})$ is a~polynomial ring of $(n-1)$ variables, which can be
identified with
the non-equivariant $K$-homology ring
$K_*(\Gr_{\SL_n})$ studied by
Lam, Schilling, and Shimozono.
If we identify the generators
of the polynomial ring with the complete symmetric functions
$h_1,\dots,h_{n-1}$, the Schubert structure sheaves are
given by the $K$-theoretic closed $k$-Schur functions. See~\cite{IIN} for a~more precise correspondence.



\subsection[Centralizer family Z]{Centralizer family $\boldsymbol{\tilde{\mathscr{Z}}}$}
\label{ssec:tZ}
We construct the solutions of~\eqref{eq:rel_Toda} from a~centralizer family $\mathscr{Z}$ defined below.
Consider the matrix equation
\begin{equation}
\label{eq:def-cent}
[A,Z]=0, \qquad
A=\begin{pmatrix}
{\rm e}^{-a_1} & -1 & 0 & \cdots &0\\
 & {\rm e}^{-a_2} & -1 & \ddots &\vdots\\
 & & \ddots & \ddots & 0\\
 & & &\ddots & -1\\
 & & & & {\rm e}^{-a_n}
\end{pmatrix},
\end{equation}
where $Z$ is an upper triangular
matrix with the entries $z_{ij}$ for $1\le i\le j\le n$. We assume $z_{11}\cdots z_{nn}\ne 0$.
$[A,Z]=0$ is equivalent to the equations
\begin{equation}
(b_i-b_j){z}_{ij}={{z}_{i,j-1}-{z}_{i+1,j}}\qquad \text{for} \quad1\le i\le j\le n. \label{eq:cent}
\end{equation}

Let us denote the affine variety over $T$ defined by these equations by $\tilde{\mathscr{Z}}$. So we have
\begin{gather*}
\mathscr{O}\bigl(\tilde{\mathscr{Z}}\bigr)=
\Rep\big[z_{ii}^{\pm 1} \ (1\le i\le n),\,
z_{ij} \ (1\le i< j\le n)
\big]/I,\\
I=\langle
(b_i-b_j)z_{ij}-z_{i,j-1}+z_{i+1,j}
\mid 1\le i< j\le n
\rangle.
\end{gather*}
There is a~natural
$\C^\times$-action
on $\tilde{\mathscr{Z}}$ by scalar
multiplication $z_{ij}\mapsto
cz_{ij}$ $(c\in \C^\times)$.
The variety $\mathscr{Z}$ is defined to be
$\tilde{\mathscr{Z}}/\C^\times$.
Thus the coordinate ring
$\mathscr{O}({\mathscr{Z}})$
is the $\Rep$-subalgebra of $\mathscr{O}\bigl(\tilde{\mathscr{Z}}\bigr)$ generated by
$z_{ij}/z_{11}$
$(1\le i\le j\le n)$.
$\mathscr{Z}$ is a~closed subscheme of
$T\times B^\vee$, where
$B^\vee$ is the
Borel subgroup of $\PGL_n(\C)$.


\begin{Remark}
In~\cite{ISY},
$\mathscr{Z}$ and
$\tilde{\mathscr{Z}}$
are denoted by $\mathscr{Z}_{\PGL_n}$
and $\mathscr{Z}_{\GL_n}$
respectively.
\end{Remark}


Let $\Rep^\reg$ be the localization of $\Rep$ by
the multiplicative set generated by $1-{\rm e}^{a_i-a_j}$ $(i\ne j)$.
Set \smash{$\mathscr{O}\bigl(\tilde{\mathscr{Z}}\bigr)^\reg
:=\Rep^\reg\otimes_{\Rep}\mathscr{O}\bigl(\tilde{\mathscr{Z}}\bigr)$}.
We often use the following.
\begin{Proposition}\label{prop:OZ_Delta}
$\mathscr{O}\bigl(\tilde{\mathscr{Z}}\bigr)^\reg$
is generated by $z_{ii}$ $(1\le i\le n)$ as
an~$\Rep^\reg$-algebra.
\end{Proposition}

\subsection[tau-functions]{$\boldsymbol{\tau}$-functions}\label{sec:tau}
We will explain that a~Zariski open subset $\mathscr{Z}^\circ$ of $\mathscr{Z}$ is isomorphic to
$\mathscr{Y}^\circ$, the compliment of the divisors defined by $Q_i$ $(1\le i\le n-1)$, whereas
$\mathscr{Z}^\circ$ is the complement of the divisor defined by the so-called $\tau$-\emph{functions}.

Let \smash{${\rm e}^{(m)}_i:=e_i({\rm e}^{-a_1},{\rm e}^{-a_2},\dots,{\rm e}^{-a_m})$} denote the $i$-th elementary symmetric polynomial in $m$ variables ${\rm e}^{-a_1},{\rm e}^{-a_2},\dots,{\rm e}^{-a_m}$.
The companion matrix $C_A$ of $A$ is described as
\begin{equation}
\label{eq:C_A}
C_A:=P^{-1}AP
=\begin{pmatrix}
0& -1 \\
&0& -1 \\
& & \ddots&\ddots \\
&& & 0&-1\\
{\rm e}^{(n)}_n& {\rm e}^{(n)}_{n-1}&\cdots &{\rm e}^{(n)}_2&{\rm e}^{(n)}_1\\
\end{pmatrix},
\end{equation}
where
\begin{equation}
\label{eq:P}
P=
\begin{pmatrix}
1\\
{\rm e}^{(1)}_1 & 1\\[1mm]
{\rm e}^{(2)}_2 & {\rm e}^{(2)}_1 & 1\\[1mm]
{\rm e}^{(3)}_3 & {\rm e}^{(3)}_2 & {\rm e}^{(3)}_1&1\\[1mm]
\vdots & \vdots & \vdots & & \ddots\\
{\rm e}^{(n-1)}_{n-1}& {\rm e}^{(n-1)}_{n-2}&\cdots &\cdots&{\rm e}^{(n-1)}_1 & 1
\end{pmatrix}.
\end{equation}
For later use, we record the following formula:
\[
 \bigl(P^{-1}\bigr)_{ij}=(-1)^{i-j}h_{i-j}({\rm e}^{-a_1},\dots,{\rm e}^{-a_{j}}).
\]

Let $I$, $J$ be subsets of $[1,n]:=\{1,\dots,n\}$,
with $I=\{i_1,\dots,i_p\}$ and $J=\{j_1,\dots,j_q\}$.
Let
$X^{J}_{I}$ denote the submatrix of $X$ consisting of its $i_1,\dots,i_p$-th rows and $j_1,\dots,j_q$-th columns.
\begin{Definition}
For $1\le i\le n$, define $\tau_i, \sigma_i\in \O\bigl(\tilde{\mathscr{Z}}\bigr)$ by
\begin{gather}\label{eq:def_of_tau}
\tau_i=
\det(ZAP)^{[1,i]}_{[1,i]},\\
\label{eq:def_of_sigma_i}
\sigma_i
=
\det(ZP)^{[1,i]}_{[1,i]},
\end{gather}
and $\tau_0=\sigma_0=1$.
\end{Definition}
In particular, we have
$
\sigma_n=\det(Z)=z_{11}\cdots z_{nn}$, $
\tau_n={\rm e}^{a_1+\dots+a_n}\sigma_n$.
The open set $\tilde{\mathscr{Z}}^\circ$
of $\tilde{\mathscr{Z}}$ is
the complement of the closed set defined by
$\tau_i= 0$, $\sigma_i=0$ ($1\le i\le n-1$).




We often use the following results on
the determinants.
%without explicitly mentioning it.
\begin{Proposition}\label{prop:Noumi}
Let $A$, $B$ be square matrices of size $n$, and $1\le i\le n$.
If $A$ is lower triangular or $B$ is upper triangular, then
\[
\det(AB)^{[1,i]}_{[1,i]}=\det A^{[1,i]}_{[1,i]}
\det B^{[1,i]}_{[1,i]}.
\]
% If,
% \[
% \det(AB)^{[1,i]}_{[1,i]}=b_{11}\cdots b_{ii}\det A^{[1,i]}_{[1,i]}.
% \]
\end{Proposition}


\begin{Proposition}\label{prop:det AB partial}
Let $A$, $B$ be square matrices of size $n$ such that $B$ is invertible, and $1\le i\le n-1$.
\[
\det(AB)^{[i+1,n]}_{[i+1,n]}
=\det\left(
\begin{matrix}
\bigl(B^{-1}\bigr)^{[1,n]}_{[1,i]}\\\hline
A^{[1,n]}_{[i+1,n]}
\end{matrix}
\right)\cdot \det B.
\]
\end{Proposition}


\subsection[Construction of Phi\_n]{Construction of $\boldsymbol{\Phi_n}$}
\begin{Theorem}[{\cite[Section~3]{IIM}}]By the map defined by~\eqref{eq:Phi}, we have
an isomorphism
$\mathscr{Y}^\circ\cong \mathscr{Z}^\circ$
of varieties over $T$.
\end{Theorem}




An element of $\tilde{\mathscr{Z}}$
is an algebraic family $\{Z_t\}$ of invertible
upper triangular matrices
parame\-trized by $t\in T$ such that $Z_t$ commute with $A_t$.
We denote the family $\{Z_t\}$ simply by $Z$.
If we assume~${Z\in \tilde{\mathscr{Z}}^\circ}$,
there exists
an upper triangular matrix $R$
and a~unipotent lower triangular matrix~$U$ (see \cite[Proposition 3.2]{IIM}), both defined as families over $T$, such that \begin{equation}\label{eq:Gauss_decomp}
P^{-1}ZAP=U^{-1}R.
\end{equation}
Because the left-hand side of~\eqref{eq:Gauss_decomp} commutes with $C_A$, we have
\begin{equation}\label{eq:AZ_to_L} UC_AU^{-1}=RC_AR^{-1}.
\end{equation}
Let $L$ denote the matrix~\eqref{eq:AZ_to_L}.
If we replace $Z$ by $cZ$ with $c\in \C^\times$,
then $R$ becomes $cR$ and we obtain the same matrix $L$.
We apply the Gauss decomposition to
$L$ as $L=MN^{-1}$. Then $M$, $N$ are matrices of the forms given in~\eqref{eq:def_of_M_and_Q} for unique $z_1,\dots,z_n, Q_1,\dots,Q_{n-1}$ (see~\cite[Section~3.2]{IIM}).
Thus we obtain
functions
$z_i$, $Q_i$
on $\mathscr{Z}^\circ=\tilde{\mathscr{Z}}^\circ/\C^{\times}$.
In view of~\eqref{eq:AZ_to_L},
we have~${\det(\zeta E-L)
=\det(\zeta E-C_A)}$, which means that
the regular functions $z_1,\dots,z_n$, $Q_1,\dots,\allowbreak Q_{n-1}$
satisfy the definiting equation for $\mathscr{Y}$.
We can also check $Q_i\ne 0$
for $1\le i\le n-1$ (see \cite[Section~3.5]{IIM}).
In this way, we have a~ring homomorphism
$\Phi_n\colon \mathscr{O}(\mathscr{Z}^\circ)\rightarrow\mathscr{O}(\mathscr{Y}^\circ)$, which is naturally a~homomorphism of $\Rep$-algebras.

The explicit form of $\Phi_n$
is determined as follows.
By abuse of notation
we simply denote $\Phi_n(z_i)$, $\Phi_n(Q_i)$ by $z_i$, $Q_i$.
Let $r_{ij}$ be the $(i,j)$-th entry of $R$.
Then, from~\eqref{eq:Gauss_decomp} and the Cauchy--Binet formula, we obtain
\begin{equation}\label{eq:r_to_minor}
\begin{aligned}
r_{11}r_{22}\cdots r_{ii}
&=
\det R^{[1,i]}_{[1,i]}=
\det\bigl(UP^{-1}ZAP\bigr)^{[1,i]}_{[1,i]}
=\det(ZAP)^{[1,i]}_{[1,i]}=\tau_i,
\end{aligned}
\end{equation}
which implies $r_{ii}=\tau_{i}/\tau_{i-1}$.
Comparing the $(i+1,i)$-th entries on both sides of
$
NM^{-1}=RC_A^{-1}R^{-1}
$, which is derived from~\eqref{eq:AZ_to_L}, we obtain
\begin{equation}\label{eq:Q_to_tau}
Q_i=\frac{r_{i+1,i+1}}{r_{ii}}=\frac{\tau_{i+1}\tau_{i-1}}{\tau_i^2}.
\end{equation}


On the other hand, from~\eqref{eq:def_of_M_and_Q} and~\eqref{eq:AZ_to_L}, $z_i$ is expressed as
\[
z_{i}
=\frac{\det\bigl(L^{-1}\bigr)^{[1,i-1]}_{[1,i-1]}}{\det\bigl(L^{-1}\bigr)^{[1,i]}_{[1,i]}}
=\frac{\det\bigl(UC_A^{-1}U^{-1}\bigr)^{[1,i-1]}_{[1,i-1]}}{\det\bigl(UC_A^{-1}U^{-1}\bigr)^{[1,i]}_{[1,i]}}
=
\frac{\det\bigl(C_A^{-1}U^{-1}\bigr)^{[1,i-1]}_{[1,i-1]}}{\det\bigl(C_A^{-1}U^{-1}\bigr)^{[1,i]}_{[1,i]}}.
\]
As $C_A^{-1}U^{-1}
=P^{-1}ZPR^{-1}$, we have
\begin{align*}
\det\bigl(C_A^{-1}U^{-1}\bigr)^{[1,i]}_{[1,i]}
&
=\det\bigl(P^{-1}ZPR^{-1}\bigr)^{[1,i]}_{[1,i]}
=\frac{\det\bigl(P^{-1}ZP\bigr)^{[1,i]}_{[1,i]}}{r_{11}\cdots r_{ii}}
=\frac{\det(ZP)^{[1,i]}_{[1,i]}}{r_{11}\cdots r_{ii}}
=\frac{\sigma_i}{\tau_i},
\end{align*}
which gives
\begin{equation}\label{eq:z_to_tau}
z_i
=\frac{\tau_i\sigma_{i-1}}{\tau_{i-1}\sigma_i}.
\end{equation}


We can construct the inverse of
$\Phi_n$ (see \cite[Section~3.4]{IIM}).
Here we provide an expression for~$\Phi_n^{-1}$.
From~\eqref{eq:Gauss_decomp}, we have
$
P^{-1}ZP=U^{-1}R C_A^{-1}$.
The entries of $U^{-1}$ can be expressed as
polynomials in $z_i$, $Q_i$, while the entries of $R$ can be written as Laurent polynomials in $z_i$, $Q_i$.
This gives rise to an expression
for the entries of $Z$ as Laurent polynomials in $z_i$, $Q_i$, which are considered as elements in $\mathscr{O}(\mathscr{Y}^\circ)$.
\begin{Example} For $n=2$, we have
\[%\label{eq:Z2}
Z=
\frac{1}{z_1z_2Q_1}
\begin{pmatrix}
z_2-{\rm e}^{-a_1}&1\\
0&z_2-{\rm e}^{-a_2}
\end{pmatrix}.
\]
For $n=3$, we have
\begin{align*}
&z_{11}=(z_1z_2z_3Q_1Q_2)^{-1}\bigl(z_2z_3-{\rm e}^{-a_1}(z_2(1-Q_2)+z_3)+{\rm e}^{-2a_1}\bigr),\\
&z_{12}=(z_1z_2z_3Q_1Q_2)^{-1}\bigl(z_2(1-Q_2)+z_3-{\rm e}^{-a_1}-{\rm e}^{-a_2}\bigr),\\
&z_{13}=(z_1z_2z_3Q_1Q_2)^{-1}.
\end{align*}
The other entries are determined by
$z_{i+1,j+1}=\omega(z_{ij})$.
\end{Example}

\begin{Remark}
As a~complex manifold, $\mathscr{Z}^\circ$ is parameterized as
\begin{equation}\label{eq:t_to_Z}
Z=\exp\bigl(
At_1+A^2t_2+\cdots+A^{n-1}t_{n-1}
\bigr)\in \mathscr{Z}^\circ
\end{equation}
by the complex parameters $t_1,t_2,\dots,t_{n-1}$.
From~\eqref{eq:t_to_Z}, we deduce the differential equation of motion $\frac{\partial }{\partial t_i}Z=A^iZ$.
Then, computing the differential $\frac{\partial}{\partial t_i}\bigl(P^{-1}ZAP\bigr)$, we obtain
\[
\frac{\partial}{\partial t_i}\bigl(P^{-1}ZAP\bigr)=P^{-1}\left(\frac{\partial}{\partial t_i}Z\right)AP=P^{-1}A^iZAP=
C_A^iU^{-1}R,
\]
where, for the last equality, we used~\eqref{eq:Gauss_decomp}.
On the other hand, by using~\eqref{eq:Gauss_decomp} again, we also obtain
\[
\frac{\partial}{\partial t_i}\bigl(P^{-1}ZAP\bigr)
=
\frac{\partial}{\partial t_i}\bigl(U^{-1}R\bigr)
=U^{-1}\left(\frac{\partial}{\partial t_i}R\right)-U^{-1}\left(\frac{\partial}{\partial t_i}U\right)U^{-1}R.
\]
Comparing these equations, we have
\[%\label{eq:der_of_Toda}
L^i=UC_A^iU^{-1}=
\left(\frac{\partial}{\partial t_i}R\right)R^{-1}-\left(\frac{\partial}{\partial t_i}U\right)U^{-1}.
\]
Since $\bigl(\frac{\partial}{\partial t_i}R\bigr)R^{-1}$ is upper triangular, we obtain \smash{$\bigl(L^i\bigr)_{<}=-\bigl(\frac{\partial}{\partial t_i}U_{t_i}\bigr)U^{-1}$} by~\eqref{eq:AZ_to_L}. Then, we have the relativistic Toda lattice~\eqref{eq:rel_Toda} as follows:
\begin{align*}%\label{eq:rel_Toda_shifted}
\frac{\partial}{\partial t_i}L
&
=
\frac{\partial}{\partial t_i}\bigl(UC_AU^{-1}\bigr)=
\left(\frac{\partial}{\partial t_i}U\right)C_AU^{-1}-UC_AU^{-1}
\left(\frac{\partial}{\partial t_i}U\right)U^{-1}
\\
&=
\left(\frac{\partial}{\partial t_i}U\right)U^{-1}L-L
\left(\frac{\partial}{\partial t_i}U\right)U^{-1}=\big[L,\bigl(L^i\bigr)_{<}\big].
\end{align*}
\end{Remark}




\section[K-theoretic double k-Schur functions]{$\boldsymbol{K}$-theoretic double $\boldsymbol{k}$-Schur functions}\label{sec:Kk}
In order to study
the map $\Phi_n$
in more detail,
we use a~recent
work~\cite{ISY} on a~symmetric function realization
of $K_*^T(\Gr_{\SL_n})$.
\subsection[Definition of K-theoretic double k-Schur functions]{Definition of $\boldsymbol{K}$-theoretic double $\boldsymbol{k}$-Schur functions}
\label{sec:def_K-k}
The simple roots of $G=SL_n(\C)$ are given by
$\alpha_i=a_i-a_{i+1}$ for $i\in I$, where
$I=\{1,\dots,n-1\}$.
Let $\theta=a_1-a_n$ denote the highest root, and let $\theta^\vee$ be its corresponding coroot.
We define an action of $\tilde{W}_G$ on $R(T)$, referred to as the \emph{level zero} action. The coroot lattice $Q^\lor$ acts by the identity, while
$W_G$ acts naturally. In particular, $s_0$ acts as $s_\theta$ where $\theta$ is the highest root.

We will work in the \emph{level zero} affine setting so we set $\alpha_0:=-\theta$.
Let $\tilde{W}_G$ denote the corresponding affine Weyl group, with the standard generators $s_0,s_1,\dots,s_{n-1}$.
Let $\Rep^\reg$ denote the localization of $\Rep$ by
the multiplicative set generated by $1-{\rm e}^{\alpha}$ where
$\alpha$ are any roots.
The \emph{twisted group algebra} $\Rep^\reg\big[\tilde{W}_G\big]$
is $\bigoplus_{w\in\tilde{W}_G }\Rep^\reg w$ with product defined by
\[
(f_1w_1)(f_2w_2)=(f_1w_1(f_2))(w_1w_2)
\qquad \text{for}\quad f_1,f_2\in\Rep^\reg,\quad w_1,w_2\in \tilde{W}_G.
\]
with the level zero action of
$\tilde{W}_G$ on $\Rep^\reg$.


For $i\in \tilde{I}=I\cup\{0\}$, define the \emph{ Demazure operators}
\begin{equation}
\label{eq:def_Demazure}
T_i=(1-{\rm e}^{\alpha_i})^{-1}(s_i-1),\qquad
D_i=T_i+1,
\end{equation}
which are considered as
elements of \smash{$\Rep^\reg\big[\tilde{W}_G\big]$}.


$T_i$ satisfies $T_i^2 = -T_i$ and the braid relations of type \smash{$A_{n-1}^{(1)}$}. Similarly, $D_i$ satisfies $D_i^2 = D_i$ and the braid relation of type \smash{$A_{n-1}^{(1)}$}. Then, for any $ w \in \hat{W}_G$ written as $w = s_{i_1} \cdots s_{i_l}$, the product
$
%\label{eq:Tw}
T_w = T_{i_1} \cdots T_{i_l}
$ and the product
$
%\label{eq:Dw}
D_w = D_{i_1} \cdots D_{i_l}
$ depend only on $w$ (see \cite{LSS:K}, see also~\cite{ISY}).

The $K$-theoretic nil-Hecke algebra $\mathbb{K}_G$ is defined to be the left $\Rep$-module
generated by~$D_w$ for $w\in \hat{W}_G$
(see \cite{LSS:K}, see also~\cite{ISY}).

We use notation
\begin{equation}
\label{eq:b_i}
b_i:=1-{\rm e}^{-a_i},\qquad 1\le i\le n.
\end{equation}

We consider the
ring \smash{$\hat{\Lambda}^{\Rep}$} of symmetric
formal power series
in the infinitely many variables~${y=(y_1,y_2,\dots)}$ with coefficients in
$\Rep$.
We can define an action of Demazure operators~$D_i$ $(0\le i\le n-1)$
on $\hat{\Lambda}^{\Rep}$ (see~\cite{ISY}).
For $1\le i\le n-1$, we
let $s_i$ act by exchanging
$a_i$ and~$a_{i+1}$
in the coefficients,
and $s_0$ by
the formula
\begin{equation}
\label{eq:s0_Omega}
s_0(f)=\frac{\Omega(b_1|y)}{\Omega(b_n|y)}s_\theta(f),\qquad
\Omega(b_i|y):=\frac{1}{\prod_{j=1}^{\infty}(1-b_i y_j)},
\end{equation}
where $s_\theta $ is the reflection with respect to $\theta$,
exchanging $a_1$ and $a_n$.

The finite Weyl group $W=\langle s_1,\dots,s_{n-1}\rangle$ is the symmetric group $S_n$.
Let \smash{$\hat{W}_G^0$} denote the set of minimal-length coset representatives of \smash{$\hat{W}_G/W_G$}. The set \smash{$\hat{W}_G^0$} naturally indexes the set of the Schubert classes of \smash{$K_*^T(\Gr_{G})$},
and we refer to an element of \smash{$\hat{W}_G^0$} as an affine Grassmann element.
For each
\smash{$x\in \hat{W}_G^0$}
\cite{ISY},
the \emph{double $K$-theoretic $k$-Schur functions}
are defined
as
\begin{equation}
\label{eq:def:tg,g}
\tilde{g}_x^{(k)}(y|b):=D_x (1),\qquad
g_x^{(k)}(y|b):=T_x (1).
\end{equation}
The $\Rep$-span
\[
\hat{\Lambda}^{\Rep}_{(n)}:=\bigoplus_{x\in \hat{W}_G^0}
\Rep \tilde{g}_x^{(k)}(y|b)
\]
is an $\Rep$-subalgebra of
\smash{$\hat{\Lambda}
_\Rep$} which is
isomorphic to
\smash{$K_*^T(\Gr_{\SL_n})$} and the Schubert structure~sheaf $\mathscr{O}_x$
is identified with \smash{$\tilde{g}_x^{(k)}(x|b)$}.
Let $\omega\in S_n$ be the cyclic permutation sending $i$ to $i+1$ with $n+1=1$ by convention.



\begin{Theorem}[\cite{ISY}]
\label{thm:ISY}
There are
isomorphisms $\beta$, $\kappa$ of $R(T)$-algebras
\[%\label{eq:main_isom}
\mathscr{O}(\mathscr{Z})
\overset{\beta}{\longrightarrow }
 \hat{\Lambda}_{(n)}^\Rep
 \overset{\kappa}{\longrightarrow}
 K_*^T(\Gr_{\SL_n}),
\]
such that
\begin{gather}
\beta(z_{ij}/z_{11})= {\rm e}^{a_i+\cdots+a_{j-1}}\Omega(b_i|y)\Omega(b_1|y)^{-1}g_{\rho_{j-i}}\big(y|\omega^i(b)\big),\nonumber
\\
\kappa(\tilde{g}_{x}^{(k)}(y|b))
= \mathscr{O}_x,
\qquad
\kappa(g_{x}^{(k)}(y|b))=\mathscr{I}_{x},\label{eq:zij_in_KkSchur}
\end{gather}
where $\rho_l=s_{l-1}\cdots s_{1}s_{0}$ for $1\le l\le n-1 $.
\end{Theorem}

\begin{Remark}
In~\cite{ISY}, it is proven that
each of the three $R(T)$-algebras has a~natural $\Rep$-Hopf-algebra structure.
\end{Remark}

\begin{Remark}
It is natural to define a~map
\smash{$\beta\colon \mathscr{O}\bigl(\tilde{\mathscr{Z}}\bigr)
\rightarrow \hat{\Lambda}^\Rep$}. We have
\begin{equation}
\beta(z_{ij})= {\rm e}^{a_i+\cdots+a_{j-1}}\Omega(b_i|y)g_{\rho_{j-i}}\big(y|\omega^i(b)\big),
\label{eq:zij_in_KkSchur_GL}
\end{equation}
where $T$ is the maximal torus of $\GL_n(\C)$. See \cite[Theorem 1.2]{ISY} for more details.
\end{Remark}

\subsection{Preliminaries on Demazure operators}
We collect some properties of $T_i$ and $D_i$ which are used in Section~\ref{sec:det_formula}.
Let $\iota$ be the involution on \smash{$\hat{\Lambda}^{R(T)}$} such that
\[
\iota(h_i(y))=\sum_{r=0}^{i-1}\binom{i-1}{r}e_{r+1}(y),\qquad
\iota({\rm e}^{a_i})={\rm e}^{-a_{n-i+1}}.
\]
\begin{Proposition}\label{prop:iota Dem}
For $i\in I$, $\iota\circ T_i=T_{n-i}\circ \iota$, $
\iota\circ D_i=D_{n-i}\circ \iota$.
\end{Proposition}
\begin{proof}
Straightforward.
\end{proof}


Let $\omega_k$ be the group automorphism of $\hat{W}_G$ such that
\begin{equation}
\label{eq:omega_k}
\omega_k(s_i)=s_{-i},
\end{equation}
with indices taken mod $n\Z$.
\begin{Proposition}[{\cite[Proposition 3.14]{ISY}}]
\label{prop:iota $K$-k}
We have
$
\iota\bigl(\tilde{g}_x^{(k)}\bigl(y|\omega^i b\bigr)\bigr)=
\tilde{g}^{(k)}_{x^{\omega_k}}\bigl(y|\omega^{-i}b\bigr)$.
\end{Proposition}

We define elements $T_{\theta}$ and $D_\theta$ of the twisted group algebra $\Rep^\reg\big[\tilde{W}_G\big]$:
\begin{equation}
\label{eq:Ttheta}
T_\theta:=\omega^{-i}\circ T_i \circ \omega^i=\frac{s_\theta-1}{1-{\rm e}^{-\theta}},\qquad
D_\theta:=\omega^{-i}\circ D_i \circ \omega^i=T_\theta+1.
\end{equation}

\begin{Remark}
For any $\alpha\in Q$, an element $T_\alpha$ of $\Rep^\reg\big[\tilde{W}_G\big]$ is defined. See \cite[Section~2.2.4]{ISY} for more details.
\end{Remark}





\begin{Lemma}\label{lem:T_and_D}
For $i\in I$,
$T_{i}\circ {\rm e}^{-a_{i+1}}={\rm e}^{-a_{i}}\circ D_{i}$, and
$T_{\theta} \circ {\rm e}^{-a_1}={\rm e}^{-a_n}\circ D_{\theta}$.
\end{Lemma}
\begin{proof}
The lemma is shown by direct calculations:
$T_{j}({\rm e}^{-a_{j+1}}f)=T_{j}({\rm e}^{-a_{j+1}})f+s_{j}({\rm e}^{-a_{j+1}})T_{j}(f)\allowbreak
={\rm e}^{-a_{j}}f+{\rm e}^{-a_{j}}T_{j}(f)={\rm e}^{-a_{j}}D_{j}(f)$.
\end{proof}


\begin{Proposition}\label{prop:T_0}
 We have
 \[
 T_0=\Omega(b_1|y)\circ T_{\theta}\circ \Omega(b_1|y)^{-1}
 =\Omega(b_n|y)^{-1}\circ T_{\theta}\circ \Omega(b_n|y).
 \]
\end{Proposition}
\begin{proof}
The first equality is shown as follows:
\begin{gather*}
\Omega(b_1|y)\circ T_{\theta}\circ \Omega(b_1|y)^{-1}(f)\\
\qquad
=
\Omega(b_1|y)
T_{\theta}\left(
\frac{f}{\Omega(b_1|y)}
\right)
=
\frac{1-b_n}{b_n-b_1}
\Omega(b_1|y)
\left(
\frac{f}{\Omega(b_1|y)}
-
\frac{s_\theta(f)}{\Omega(b_n|y)}
\right)\\
\qquad
=
\frac{1-b_n}{b_n-b_1}
\left(
f
-
\frac{\Omega(b_1|y)}{\Omega(b_n|y)}s_{\theta}(f)
\right)=T_0(f).
\end{gather*}
The second equality follows from the fact that $\Omega(b_1|y)\Omega(b_n|y)$ commutes with $T_{\theta}$.
\end{proof}

\begin{Lemma}\label{lem:T_on_h}
For
$1\le i\le n-1$ and $j-i\le n-1$, and $m\ge 0$,
\begin{equation}
D_{i}
h_m({\rm e}^{-a_{i+1}},{\rm e}^{-a_{i+2}},\dots,{\rm e}^{-a_j})
=
h_m({\rm e}^{-a_{i}},{\rm e}^{-a_{i+1}},{\rm e}^{-a_{i+2}},\dots,{\rm e}^{-a_j}).
\label{eq: Di hm}
\end{equation}
For $1\le j\le n-1$, and $m\ge 0$,
\begin{equation}
D_{\theta}
h_m({\rm e}^{-a_{1}},{\rm e}^{-a_{2}},\dots,{\rm e}^{-a_j})
=
h_m({\rm e}^{-a_{n}},{\rm e}^{-a_{1}},{\rm e}^{-a_{2}},\dots,{\rm e}^{-a_j}).
\label{eq: D theta hm}
\end{equation}
\end{Lemma}
\begin{proof}
Note
\begin{equation*}
D_i (1-{\rm e}^{-a_{i+1}}u)^{-1}=
(1-{\rm e}^{-a_{i}}u)^{-1}
(1-{\rm e}^{-a_{i+1}}u)^{-1}.
\end{equation*}
As $\prod_{s=i+2}^j(1-{\rm e}^{-a_s}u)^{-1}$ is invariant under $s_i$, we obtain
$D_i\prod_{s=i+1}^j(1-{\rm e}^{-a_s}u)^{-1}=
\prod_{s=i}^j(1-{\rm e}^{-a_s}u)^{-1}$, which implies~\eqref{eq: Di hm}.
The proof of~\eqref{eq: D theta hm} is given similarly.
\end{proof}


\section{Some automorphisms}\label{sec:auto}
For later use, we introduce
some automorphisms of the algebras $\mathscr{O}\bigl(\tilde{\mathscr{Z}}\bigr), \hat{\Lambda}^{\Rep}$
and study their properties.
\subsection[Automorphism sigma]{Automorphism $\boldsymbol{\sigma}$}\label{sec:sigma}

We define an automorphism $\sigma$
of $\mathscr{O}\bigl(\tilde{\mathscr{Z}}\bigr)$ as an $\Rep$-algebra
by
\begin{align}\label{eq:def_of_sigma}
\sigma(Z)&=ZA^{-1}.
\end{align}
Explicitly,
for $1\le i\le j\le n$,
we have
\[
%\label{eq:sigma_z_ij}
\sigma(z_{ij})={\rm e}^{a_j}z_{ij}+
{\rm e}^{a_{j-1}+a_j}z_{i,j-1}+\cdots+{\rm e}^{a_i+\cdots+a_j}
z_{ii}.
\]
In particular,
we have
\begin{equation}
\label{eq:sigma_zii}
\sigma(z_{ii})={\rm e}^{a_i}z_{ii}.
\end{equation}
From~\eqref{eq:def_of_tau} and~\eqref{eq:def_of_sigma_i}, we have
\begin{equation}\label{eq:sigma_tau_i}
\sigma(\tau_i)=\sigma_i.
\end{equation}


Via the isomorphism $\beta$
we consider the corresponding automorphism on \smash{$\hat{\Lambda}_{(n)}^\Rep$}
and denote it also by $\sigma$.
For example,
we have
\begin{equation}
\sigma(\Omega(b_i|y))={\rm e}^{a_i}\Omega(b_i|y)\label{eq:sigma(Omega)}.
\end{equation}
In the next section, we prove
\begin{equation}
\sigma(\ks_{\rho_i}^{(k)}(y|b))={\rm e}^{a_i-a_n}\kst_{\rho_i}^{(k)}(y|b).\label{eq:sigma(tg)}
 \end{equation}
Since we know
$g_{\rho_i}^{(k)}(y|0)=h_i(y)$,
we have
$
\sigma(h_i(y))=1+h_1(y)+\cdots+h_i(y)$.
In fact, $\sigma$
on~\smash{$\hat{\Lambda}^\Rep$} can be defined
as the $\Rep$-linear map
sending $f(y_1,y_2,\dots)\in \hat{\Lambda}$
to $f(1,y_1,\allowbreak y_2,\dots)$.
This automorphism
already appeared in the study of the
non-equivariant version of $K$-Peterson isomorphism in
\cite{BMS,IIN}.

\subsection[Image of g\^{}(k)\_la(y|b) under sigma for k-small la]{Image of $\boldsymbol{{g}^{(k)}_\la(y|b)}$ under $\boldsymbol{\sigma}$ for $\boldsymbol{k}$-small $\boldsymbol{\la}$}

Let $\mathrm{diag}(\lambda)$ denote the \textit{main diagonal} of $\lambda$, that is, the set of boxes at the $(i,i)$-th position for~${i=1,2,\dots}$.
For $x\in \mathrm{diag}(\lambda)$, let $r(x)$ be the $n$-residue of the box that is furthest to the right from $x$, and $b(x)$ the $n$-residue of the box that is furthest below from $x$.
\begin{Proposition}\label{prop:v1_of_kKShur}
If $\lambda$ is $k$-small, we have
\[
\sigma\bigl(g^{(k)}_\lambda(y|b)\bigr)
=
\Bigg(\prod_{x\in \mathrm{diag}(\lambda)}{\rm e}^{a_{r(x)+1}-a_{b(x)}}\Bigg)
\tilde{g}^{(k)}_\lambda(y|b).\]
\end{Proposition}

\begin{Example}
When $n=6$ and
\[
\ytableausetup{smalltableaux}\lambda=
 \begin{ytableau}
 0 & 1& 2\\ 5&0&1\\4
 \end{ytableau}
 ,
 \]
the main diagonal consists of two boxes $x_1$, $x_2$, where $x_i$ is at the $(i,i)$-th position.
Since
$r(x_1)=2$,
$b(x_1)=4$,
$r(x_2)=1$, and
$b(x_2)=0$,
we have \smash{$\sigma(g^{(5)}_{\lambda}(y|b))={\rm e}^{(a_3-a_4)+(a_2-a_6)}\Tilde{g}^{(5)}_\lambda(y|b)$}.
\end{Example}

\begin{Lemma}\label{lem:sigma T0}
$\sigma \circ T_0={\rm e}^{\theta} D_0 \circ \sigma$.
\end{Lemma}
\begin{proof}
Let \smash{$f(y|b)\in \hat{\Lambda}^\Rep$}. Then
by using~\eqref{eq:iota Omega}, we have
\begin{align*}
\sigma T_0 f(y|b)
&=\sigma
\left(
\frac{1}{1-{\rm e}^{-\theta}}(\frac{\Omega(b_1|y)}{\Omega(b_n|y)}s_\theta f(y|b) -f(y|b))
\right)\\
&=
\frac{1}{1-{\rm e}^{-\theta}}\left({\rm e}^{\theta}\frac{\Omega(b_1|y)}{\Omega(b_n|y)} \sigma(( s_\theta f)(y| b)) -(\sigma f)(y|b)\right)\\
&=
\frac{1}{1-{\rm e}^{-\theta}}\left({\rm e}^{\theta}\frac{\Omega(b_1|y)}{\Omega(b_n|y)}s_\theta(( \sigma f)(y| b))-{\rm e}^{\theta} \sigma f(y| b) +{\rm e}^{\theta}\sigma f(y| b)-\sigma f(y|b)\right)\\
&={\rm e}^{\theta}
T_0\sigma f(y|b)
+{\rm e}^{\theta}\sigma f(y|b)={\rm e}^{\theta}
D_0\sigma f(y|b),
\end{align*}
where we use that $\sigma$ commutes with the action of $s_\theta$.
\end{proof}

\begin{proof}[Proof of Proposition~\ref{prop:v1_of_kKShur}]
We use induction on the number of boxes of $\la$. The case $\la=\varnothing$ is
obvious. Suppose~${\la\ne \varnothing}$.
There is a~box removable from $\la$, with the $n$-residue say $i$. Let $\mu$ be the partition obtained from $\la$ by removing the box.
We consider the case when $i=0$.
One easily sees that~${r(x)+1}$, $b(x)\notin\{1,n\}$ for $x\in \mathrm{diag}(\mu)$. It follows that $s_\theta (e(\mu))=e(\mu)$. We also note that~${{\rm e}^\theta e(\mu)=e(\lambda)}$,
\begin{align*}
\sigma\bigl(g^{(k)}_\lambda(y|b)\bigr)
&=\sigma\bigl(T_0 g^{(k)}_\mu(y|b)\bigr)={\rm e}^{\theta}{D}_0\sigma\bigl( g^{(k)}_\mu(y|b)\bigr)\qquad \text{by Lemma~\ref{lem:sigma T0}}\\
&={\rm e}^{\theta}{D}_0\bigl(e(\mu)\tilde{g}^{(k)}_\mu(y|b)\bigr)={\rm e}^{\theta}e(\mu){D}_0\bigl(\tilde{g}^{(k)}_\mu(y|b)\bigr)
=e(\lambda)\tilde{g}^{(k)}_\lambda(y|b).
\end{align*}
The case when $i\ne 0$ is
left to the reader since it is
similar and easier.
\end{proof}

\begin{Remark}
The assumption that $\lambda$ is $k$-small in Proposition~\ref{prop:v1_of_kKShur} is mandatory.
Indeed, if we consider the case when $n=3$ and $w=s_2s_1s_0$, the associated $2$-bounded partition $\lambda=(2,1)$ is not $2$-small, and
\smash{$\sigma\bigl(g^{(2)}_w(y|b)\bigr)$} equals
\begin{gather*}
\mathcal{T}_2\bigl({\rm e}^{a_2-a_3}\otimes \tilde{g}^{(2)}_{10}(y|b)\bigr)\\
\qquad
=\bigl(1+{\rm e}^{a_3-a_2}\bigr)\otimes \tilde{g}^{(2)}_{10}(y|b)+{\rm e}^{a_3-a_2}\otimes T_2\tilde{g}^{(2)}_{10}(y|b) \\
\qquad=1\otimes \tilde{g}^{(2)}_{10}(y|b)+{\rm e}^{a_3-a_2}\otimes D_2\tilde{g}^{(2)}_{10}(y|b)
=1\otimes \tilde{g}^{(2)}_{10}(y|b)+{\rm e}^{a_3-a_2}\otimes \tilde{g}^{(2)}_{210}(y|b).
\end{gather*}
\end{Remark}







\subsection[Regular functions c\_i]{Regular functions $\boldsymbol{c_i}$}\label{sec:ci}




For arbitrary matrix $Z$ that commutes with $A$, there are unique scalars $c_i$ satisfying
\begin{equation}\label{eq:def_of_c_i}
Z=c_0E+c_1A+\cdots+c_{n-1}A^{n-1}.
\end{equation} The existence of such
scalars is assured by the fact that $A$ is conjugate to a~companion matrix.
We will consider each $c_i$ as an element of $\mathscr{O}\bigl(\tilde{\mathscr{Z}}\bigr)$.

\begin{Example} If
$n=3$ we have expressions
\[
c_0=z_{11}+{\rm e}^{-a_1}z_{12}+{\rm e}^{-a_1-a_2}z_{13},\qquad
c_1=-z_{12}-({\rm e}^{-a_1}+{\rm e}^{-a_2})z_{13},\qquad
c_2=z_{13}
\]
by comparing the $1$st row of~\eqref{eq:def_of_c_i}.
\end{Example}




By comparing the diagonal
entries, we have
\smash{$
z_{ii}=c_0+c_1 {\rm e}^{-a_i}+
\cdots+c_{n-1}{\rm e}^{-(n-1)a_i}$}.
%\label{eq:zii_c_1}
We also define \smash{$c_0^{(j)},c_1^{(j)},\dots,c_{n-1}^{(j)}\in \mathscr{O}\bigl(\tilde{\mathscr{Z}}\bigr)$}
for $j\in \Z$ by
\begin{equation}\label{eq:def_of_c_i_(j)}
A^jZ=c_0^{(j)}E+c_1^{(j)}A+\cdots+c_{n-1}^{(j)}A^{n-1}.
\end{equation}
Note that we have $c_j^{(0)}=c_j$.
Comparing the diagonal entries on both sides of~\eqref{eq:def_of_c_i_(j)}, we have
\begin{equation}
{\rm e}^{-ja_i}z_{ii}=c_0^{(j)}+c_1^{(j)} {\rm e}^{-a_i}+
\cdots+c_{n-1}^{(j)}{\rm e}^{-(n-1)a_i}.\label{eq:zii_c_2}
\end{equation}









\begin{Theorem}\label{thm:det_c_to_tausigma}
We have
\begin{equation}\label{eq:det_c_to_tausigma}
\begin{gathered}
\sigma_i=
\begin{vmatrix}
c^{(0)}_{0} & c^{(0)}_{1} & \dots & c^{(0)}_{i-1}\\
c^{(1)}_{0} & c^{(1)}_{1} & \dots & c^{(1)}_{i-1}\\
\vdots & \vdots & \dots & \vdots \\
c^{(i-1)}_{0} & c^{(i-1)}_{1} & \dots & c^{(i-1)}_{i-1}\\
\end{vmatrix}, \qquad
\tau_i=
\begin{vmatrix}
c^{(1)}_{0} & c^{(1)}_{1} & \dots &c^{(1)}_{i-1}\\
c^{(2)}_{0} & c^{(2)}_{1} & \dots &c^{(2)}_{i-1}\\
\vdots & \vdots & \dots & \vdots\\
c^{(i)}_{0} & c^{(i)}_{1} & \dots &c^{(i)}_{i-1}\\
\end{vmatrix}.
\end{gathered}
\end{equation}
\end{Theorem}

\begin{proof}
Since $P^{-1}$ is lower unitriangular, we use
Proposition~\ref{prop:Noumi} to have
\[
\begin{aligned}
\sigma_i=\det(ZP)^{[1,i]}_{[1,i]}
=
\det\bigl(P^{-1}ZP\bigr)^{[1,i]}_{[1,i]}=
\det\bigl(c_0E+c_1C_A+\dots+c_{n-1}C_A^{n-1}\bigr)^{[1,i]}_{[1,i]}.
\end{aligned}
\]
By comparing the 1-st row of
\[
C_A^{j-1}\bigl(c_0E+c_1C_A+\cdots+c_{n-1}C_A^{n-1}\bigr)
=c_0^{(j-1)}E+c_1^{(j-1)}C_A+\cdots+c_{n-1}^{(j-1)}C_A^{n-1},
\]
we see that the
$(j,l)$-th entry of $c_0E+c_1C_A+\cdots+c_{n-1}C_A^{n-1}$
is
\smash{$(-1)^{j+l}c_l^{(j-1)}$}.
Hence the first equation of~\eqref{eq:det_c_to_tausigma} holds.
The equation for $\tau_i$ follows from this by applying $\sigma^{-1}$ because of~\eqref{eq:sigma_tau_i}.\looseness=-1
\end{proof}

\subsection[Involution iota on O(Z)]{Involution $\boldsymbol{\iota}$ on $\boldsymbol{\mathscr{O}\bigl(\tilde{\mathscr{Z}}\bigr)}$}
\label{sec:iota on OZ}

Let $J$ be the permutation matrix
of the longest element $w_\circ$ of $S_n$.
Explicitly, $J$ is $\sum_{i=1}^n E_{i,n-i+1}$.
Let $\iota$ be a~ring automorphism of $\mathscr{O}\bigl(\tilde{\mathscr{Z}}\bigr)$ defined by
$\iota({\rm e}^\gamma)=
{\rm e}^{-w_\circ \gamma}$ (in particular
$\iota({\rm e}^{a_i})={\rm e}^{-a_{n-i+1}}$), and
\[%\label{eq:def_of_S}
\sum_{i=0}^{n-1}
\iota(c_i)A^{-i}
=\left(\sum_{i=0}^{n-1}
c_iA^{i}\right)^{-1}.
\]
It is not difficult to show $\iota(C_A)=JC_A^{-1}J$.
Since $C_A=P^{-1}AP$, we obtain
\[\sum_{i=0}^{n-1}
\iota(c_i)C_A^{-i}
=\left(\sum_{i=0}^{n-1}
c_iC_A^{i}\right)^{-1}
\] and
\begin{align}
\iota\bigl(P^{-1}ZP\bigr)
&=\iota\bigl(c_0+c_1C_A+\dots+c_{n-1}C_A^{n-1}\bigr) \nonumber\\
&=J\bigl(\iota(c_0)+\iota(c_1)C_A^{-1}+\dots+\iota(c_{n-1})C_A^{-(n-1)} \bigr)J\nonumber\\
&=J\bigl(c_0+c_1C_A+\dots+c_{n-1}C_A^{n-1} \bigr)^{-1}J=JP^{-1}Z^{-1}PJ.\label{eq:S_PZP}
\end{align}
Comparing the diagonal entries on both sides of~\eqref{eq:S_PZP}, we have
\begin{equation}
\label{eq:S_zii}
\iota(z_{ii})=z_{n-i+1,n-i+1}^{-1}.
\end{equation}

Identifying $z_{ii}$ with $\Omega(b_i|y)$, we have
\begin{equation}\label{eq:iota Omega}
\iota(\Omega(b_i|y))=\Omega(b_{n-i+1}|y)^{-1}.
\end{equation}


\begin{Remark}
The automorphism $\iota$ is induced by a~diagram automorphism of the affine Dynkin diagram. See \cite[Sections~2.7 and~3.3]{ISY}.
\end{Remark}

\begin{Proposition}\label{prop:S_and_sigma}
$\iota$ commutes with $\sigma$.
\end{Proposition}
\begin{proof}
Note that $\iota$ and $\sigma$ are naturally extended
as ring automorphisms
of $\mathscr{O}\bigl(\tilde{\mathscr{Z}}\bigr)^\Delta$. In view of
Proposition~\ref{prop:OZ_Delta}, we only need to check on the generators $z_{ii}$
and $f\in R(T)$.
From~\eqref{eq:S_zii} and~\eqref{eq:sigma_zii}, we have
\begin{align*}
 \sigma (\iota(z_{ii}))
 =\sigma\bigl(z_{n-i+1,n-i+1}^{-1}\bigr)=
 {\rm e}^{-a_{n-i+1}}
 z_{n-i+1,n-i+1}^{-1}=\iota({\rm e}^{a_i}z_{ii})=\iota(\sigma(z_{ii})).
\end{align*}
Since $\sigma$ is $R(T)$-linear, it is clear that
$\iota(\sigma(f))=\sigma(\iota(f))=\iota(f)
$ for $f\in R(T)$.
\end{proof}

\begin{Proposition}
\label{prop:S_on_tau}
We have
\begin{align}
\label{eq:sigma tau_i}
\iota(\tau_i)&=\frac{\tau_{n-i}}{\tau_n},\qquad
\iota(\sigma_i)=\frac{\sigma_{n-i}}{\sigma_n}.
\end{align}
\end{Proposition}
\begin{proof}
By~\eqref{eq:S_PZP}, we have
\begin{align*}
\iota(\sigma_i)
&=\iota\left(\det\bigl(P^{-1}ZP\bigr)_{[1,i]}^{[1,i]}\right)=\det\bigl(JP^{-1}Z^{-1}PJ\bigr)_{[1,i]}^{[1,i]}=\det\bigl(P^{-1}Z^{-1}P\bigr)_{[n-i+1,n]}^{[n-i+1,n]}\\
&=
\det\bigl(P^{-1}Z^{-1}P\bigr)\cdot\det\bigl(P^{-1}ZP\bigr)_{[1,n-i]}^{[1,n-i]}=
\det(Z^{-1})\cdot\det\bigl(P^{-1}ZP\bigr)_{[1,n-i]}^{[1,n-i]}\\
&=\frac{\sigma_{n-i}}{\Omega(b_1|y)\cdots \Omega(b_n|y)}
=\frac{\sigma_{n-i}}{\sigma_n},
\end{align*}
where for the fourth equality
we use
the fact
\smash{$\det A^{[m+1,n]}_{[m+1,n]}
=\det A \cdot\det\bigl(A^{-1}\bigr)^{[1,m]}_{[1,m]}$},
which holds for any invertible $A$.
The first equality of~\eqref{eq:sigma tau_i} is obtained from Proposition~\ref{prop:S_and_sigma} and~\eqref{eq:sigma_tau_i}.
\end{proof}


\subsection[Basic properties of Phi\_n]{Basic properties of $\boldsymbol{\Phi_n}$}
We collect some basic properties of $\Phi_n$ which will be used below.

\begin{Proposition}\label{prop:S_commutes}
We have
\begin{align}
 \iota\circ \Phi_n&=\Phi_n \circ \iota.\label{eq:S_Phi}
\end{align}
\end{Proposition}
\begin{proof}
 For~\eqref{eq:S_Phi}, it suffices to show on generators $z_i$, and $Q_i$.
From Proposition~\ref{prop:S_on_tau}, we have
\[
(\iota\circ \Phi_n)(z_i)
=\iota\left(\frac{\tau_i\sigma_{i-1}}{\tau_{i-1}\sigma_i}\right)
=\frac{\tau_{n-i}\sigma_{n-i+1}}{\sigma_{n-i}\tau_{n-i+1}}=\Phi_n(z_{n-i+1}^{-1})
=(\Phi_n\circ \iota)(z_i)\]
and
\[
(\iota\circ\Phi_n)(Q_i)=\iota\left(\frac{\tau_{i-1}\tau_{i+1}}{\tau^2_{i}}\right)=
\frac{\tau_{n-i+1}\tau_{n-i-1}}{\tau^2_{n-i}}=\Phi_n(Q_{n-i})
=(\Phi_n\circ \iota)(Q_i),
\]
which concludes the proposition.
\end{proof}

Let us denote
the characteristic polynomial of \smash{$L_{[1,i]}^{[1,i]}$}
by
\smash{$
\chi_i(\zeta):=\det(\zeta E-L)_{[1,i]}^{[1,i]}$}.
Then we~have from Appendix~\ref{sec:F}
\[
\chi_i(\zeta)=
\zeta^i-F^{(i)}_1\zeta^{i-1}+\dots+(-1)^{i}F^{(i)}_i,
\]
where
\begin{equation}
\label{eq:Fij}
F_m^{(i)}:=\sum_{\substack{J\subset [i]\\ |J|=m}}
\prod_{\substack{j\in J,\\j+1\notin J}}
(1-Q_j)\prod_{j\in J}z_j,
\end{equation}
for
$0\le m\le i\le n$, with $F_0^{(i)}=1$.

\begin{Lemma}\label{lem:SF_to_c}
For $1\le m\le i\le n$, we have
\[
(\iota\circ\Phi_n )\bigl(F^{(i)}_m\bigr)=\frac{(-1)^{m}}{\tau_{n-i}}
\begin{vmatrix}
c^{(1)}_{0} & c^{(1)}_{1} & \dots & c^{(1)}_{n-i-2}&c^{(1)}_{n-i+m-1}\\
c^{(2)}_{0} & c^{(2)}_{1} & \dots & c^{(2)}_{n-i-2}&c^{(2)}_{n-i+m-1}\\
\vdots & \vdots & \dots & \vdots &\vdots\\
c^{(n-i)}_{0} & c^{(n-i)}_{1} & \dots & c^{(n-i)}_{n-i-2}&c^{(n-i)}_{n-i+m-1}\\
\end{vmatrix}.
\]
\end{Lemma}
 \begin{proof}
We first compute $\Phi_n(\chi_i(\zeta))$ as follows
\begin{gather*}
\det\bigl(U^{-1}\bigr)_{[1,i]}^{[1,i]}\cdot
\Phi_n(\chi_i(\zeta))
\cdot\det R_{[1,i]}^{[1,i]}\\
\qquad=
\det\bigl(U^{-1}\bigr)_{[1,i]}^{[1,i]}\cdot
\det(\zeta E-\Phi_n(L))_{[1,i]}^{[1,i]}
\cdot\det R_{[1,i]}^{[1,i]}\\
\qquad=
\det\bigl( U^{-1}(\zeta E-\Phi_n(L))R\bigr)_{[1,i]}^{[1,i]}\qquad\text{by Proposition~\ref{prop:Noumi}}
\\
\qquad=
\det\bigl((\zeta E-C_A) U^{-1}R\bigr)_{[1,i]}^{[1,i]}
\qquad\text{since $\Phi_n(L)=UC_AU^{-1}$ from
\eqref{eq:AZ_to_L}}\\
\qquad={\det\bigl((\zeta E-C_A)P^{-1}ZAP\bigr)_{[1,i]}^{[1,i]}} \qquad\text{by~\eqref{eq:Gauss_decomp}}.
\end{gather*}
Since \smash{$\det\bigl(U^{-1}\bigr)_{[1,i]}^{[1,i]}=1$} and
\smash{$\det R_{[1,i]}^{[1,i]}=\tau_i$} (see~\eqref{eq:r_to_minor}), we deduce that
\begin{equation}\label{eq:PhiL_Z}
\Phi_n(\chi_i(\zeta))
={\tau_i}^{-1}{\det\bigl((\zeta E-C_A)P^{-1}ZAP\bigr)_{[1,i]}^{[1,i]}}.
\end{equation}
Let us apply $\iota$ on both sides of~\eqref{eq:PhiL_Z}.
By using Proposition~\ref{prop:S_commutes}, we can verify
\[
\iota\bigl(P^{-1}ZAP\bigr)=
JP^{-1}Z^{-1}A^{-1}PJ.
\]
Using this together with $\iota(C_A)=JC_A^{-1}J$, we have
\begin{align*}
(\iota\circ \Phi_n)(\chi_i(\zeta))
&=(\tau_n/\tau_{n-i})\cdot{\det\bigl(J(\zeta E-C_A^{-1})P^{-1}Z^{-1}A^{-1}PJ\bigr)_{[1,i]}^{[1,i]}}\\
&=(\tau_n/\tau_{n-i})\cdot\det\bigl(\bigl(\zeta E-C_A^{-1}\bigr)P^{-1}Z^{-1}A^{-1}P\bigr)_{[n-i+1,n]}^{[n-i+1,n]}\\
&=\frac{1}{\tau_{n-i}}
\det\left(
\begin{matrix}
\bigl(P^{-1}AZP\bigr)_{[1,n-i]}^{[1,n]}\\
\hline
\bigl(\zeta E-C_A^{-1}\bigr)_{[n-i+1,n]}^{[1,n]}
\end{matrix}
\right),
\end{align*}
where in the
last equality we use Proposition~\ref{prop:det AB partial}, and $\det\bigl(P^{-1}AZP\bigr)=\tau_n$.
Comparing the coefficients of $\zeta^{i-m}$ on both sides of this equality, we obtain
\[
(\iota\circ\Phi_n)\bigl(F^{(i)}_m\bigr)=
\frac{\det\bigl(P^{-1}AZP\bigr)_{[1,n-i]}^{[1,n-i-1]\cup\{n-i+m\}}}{\tau_{n-i}}.
\]
Because we know $P^{-1}AZP=\bigl((-1)^{j+l}c^{(j)}_{l}\bigr)_{1\le j,l\le n}$ (see the proof of Theorem~\ref{thm:det_c_to_tausigma}), we obtain the desired result, after eliminating unnecessary signs.
\end{proof}

\section{Quantum double Grothendieck polynomials}\label{sec:Groth}
We use notation of binary operations $x\oplus y=x+y-xy$ and
$x\ominus y=(x-y)/(1-y)$.
We also denote $-x/(1-x)=0\ominus x$ by $\overline{x}$.


\subsection{Quantum double Grothendieck polynomials}
\label{sec:G^Q(z|eta)}
We recall
Lenart--Maeno's
quantum double Grothendieck polynomials.



Let us consider
the two sets of variables
$z=(z_1,\dots,z_n)$ and
$\eta=(\eta_1,\dots,\eta_n)$, and the polynomial ring
$\Z[Q][z_1,\dots,z_n,\eta_1,\dots,\eta_n]$,
where $Q_1,\dots,Q_{n-1}$ are the Novikov variables.
For $i\in I$, define
\begin{equation}
\label{eq:DiQ,TiQ}
T_i^Q
=\frac{s_i^{(\eta)}-1}{\eta_{i+1}\ominus \eta_{i}},\qquad
D_i^Q
=1+\frac{s_i^{(\eta)}-1}{\eta_{i+1}\ominus \eta_{i}},
\end{equation}
as linear endomorphisms of
$\Z[Q][z_1,\dots,z_n,\eta_1,\dots,\eta_n]$.
Define
\begin{gather}
\label{eq:G_longest}
\Gro{\wz}(z|\eta)=
\prod_{i=1}^{n-1}
\fac_i,\\
\fac_i=\sum_{j=0}^i
(-1)^j
(1-\eta_{n-i})^j
F_j^{(i)}(z,Q).
\label{eq:def phi}
\end{gather}
Note that $\fac_i$ can be written as
\begin{equation}
\label{eq:phi det}
\fac_i=\det(E-(1-\eta_{n-i})L)^{[1,i]}_{[1,i]}
\end{equation}
if $L$ is the matrix given as in Section~\ref{sec:Rel Toda}.

There exists a~unique collection of polynomials
\[
\bigl\{\Gro{w}(z|\eta)
\in \Z[Q][z_1,\dots,z_{n},\eta_1,\dots,\eta_{n}]
 \mid w\in S_n\bigr\},\]
where $\mathfrak{G}_{w_\circ}^Q$ corresponding to the longest element $w_\circ$ is given by~\eqref{eq:G_longest}, and each $\mathfrak{G}_{w}^Q$ for general~$w$ is characterized by the recursive relation
\begin{equation}
 \label{eq:Di_Gr}
D_i^Q\Gro{w}(z|\eta)=\begin{cases}
\Gro{s_iw}(z|\eta) & \text{if $s_i w<w$},\\
\Gro{w}(z|\eta) & \text{if $s_i w>w$}.
\end{cases}
\end{equation}
\begin{Remark}\label{rem:e_vs_eta}
We follow the notation in~\cite{MNS2} as closely as possible, however, there are some
unavoidable differences as outlined below.
The variables $y_i$ used for equivariant parameters in~\cite{MNS2} are denoted here by $\eta_i$.
This change is necessary because we reserve $y_i$ for the variables of symmetric functions.
We treat the variable $\eta_i$ as
an element of $\Rep$ via the correspondence
\begin{equation}
\eta_i=1-{\rm e}^{a_{n-i+1}}=\overline{b}_{n-i+1},\qquad1\le i\le n.
\label{eq:eta-b}
\end{equation}
Note that our $a_i$ corresponds to $-\epsilon_i$ in~\cite{MNS2}.
Additionally, we use $z_i=1-x_i$ with $x_i$
the variable from~\cite{MNS2}.


\end{Remark}

\begin{Proposition}For $i\in I$, we have
 \begin{equation}
D_{n-i}\circ \Phi_n=\Phi_n \circ D_i^Q,\qquad
D_{n-i}\circ \tilde{\Phi}_n=\tilde{\Phi}_n \circ D_i^Q.
\label{eq:D_and_Phi}
\end{equation}
\end{Proposition}
\begin{proof}
The operators on both sides
coincide on $z_i$, $Q_i$,
because $\tau_i$ and $\sigma_i$ are $S_n$-invariant. Since~$\Phi_n$ \big(and $\tilde{\Phi}_n$\big) is $\Rep$-linear,
it suffices to see the equality
on $\Rep$. In fact, we have
\begin{equation*}
1+\frac{s_i^{(\eta)}-1}{\eta_{i+1}\ominus \eta_{i}}
=1+
\frac{s_{n-i}^{(b)}-1}{\overline{b}_{n-i}\ominus \overline{b}_{n-i+1}}
=1+
\frac{s_{n-i}^{(b)}-1}{b_{n-i+1}\ominus b_{n-i}},
\end{equation*}
under the
identification~\eqref{eq:eta-b},
where \smash{$s_i^{(b)}$} exchanges $b_i$
and $b_{i+1}$,
and we use $\overline{x}\ominus \overline{y}=y\ominus x$.
\end{proof}


\subsection[Involution iota on the equivariant quantum K-ring]{Involution $\boldsymbol{\iota}$ on the equivariant quantum $\boldsymbol{K}$-ring}

Let $\iota$ be a~ring homomorphism of
$\Rep[Q]\big[z_1^{\pm 1},\dots,z_n^{\pm 1}\big]$
such that
\[
\iota({\rm e}^{a_i})={\rm e}^{-a_{n-i+1}},\qquad
\iota(z_{i})=z_{n-i+1}^{-1},\qquad
\iota(Q_i)=Q_{n-i}.
\]
Obviously,
$\iota$ satisfies $\iota^2=1$.
It is easy to prove that
\begin{equation}\label{eq:S_and_D_i}
 \iota\circ D_i^Q=D_{n-i}^Q\circ \iota.
 \end{equation}
 holds for $i\in I$.
\begin{Proposition}
$\iota$ preserves the ideal $\mathscr{I}_n^Q$ of $\Rep[Q]\big[z_1^{\pm 1},\dots,z_n^{\pm 1}\big]$.
\end{Proposition}
\begin{proof}
We claim that
\smash{$
z_1\cdots z_n \iota\bigl(F_i^{(n)}\bigr)=F_{n-i}^{(n)}$}, \smash{$
e_n^{(n)}\iota\bigl(e_i^{(n)}\bigr)=e_{n-i}^{(n)}
$}
for $1\le i\le n-1$.
The second identity is obvious.
For $J\subset [1,n]$, $ |J|=i$, set
$K:=\{i'\mid i\notin J\}$, with $i':=n-i+1$. Then
$z_1\cdots z_n \prod_{j\in J}\iota(z_i)=\prod_{l\in K}z_l$.
We see that $j\in J, j+1\notin J$
is equivalent to
$(j+1)'\in K$, $j'\notin K$.
Thus
\[
\prod_{j\in J, j+1\notin J}
(1-Q_j)=
\prod_{l\in K, l+1\notin K}
(1-Q_l).
\]
Hence the claim holds and we have $z_1\cdots z_n \iota\bigl(F_i^{(n)}\bigr)=F_{n-i}^{(n)}$.
Since $F_n^{(n)}=z_1\cdots z_n$, and $e_n^{(n)}={\rm e}^{-(a_1+\cdots+a_n)}=1$, we are done.
\end{proof}

The next result will be proved in Section~\ref{sec:max factor}.
\begin{Proposition}\label{prop:long_iota_invariant}
We have $\iota\bigl(\mathfrak{G}_{w_\circ}^{Q}\bigr)=\mathfrak{G}_{w_\circ}^{Q}$.
\end{Proposition}

For $w\in S_n$, let
$w^*=w_\circ w w_\circ;$
$w^*$ is obtained from any
reduced expression of $w$ by replacing~$s_i$ with $s_{n-i}$.
Note that the map $w\mapsto w^*$ preserves the Bruhat order.
\begin{Proposition}\label{prop:S_star}
Let $w\in S_n$. Then
\begin{equation}
\iota\bigl(\Gro{w}(z|\eta)\bigr)=
\Gro{w^*}(z|\eta).
\label{eq:S_star}
\end{equation}
\end{Proposition}
\begin{proof}
We use induction on $\ell(w_\circ)-\ell(w)$.
Since $\Gro{w_\circ}$
is $\iota$-invariant (see Proposition~\ref{prop:long_iota_invariant})
and $w_\circ^*=w_\circ$, \eqref{eq:S_star} holds for $w=w_\circ$.
Suppose $w\in S_n$ satisfies
$s_iw>w$ for some $1\le i\le n-1$.
Then we have $(s_iw)^*>w^*$.
It follows that $s_{n-i}(s_iw)^*<(s_iw)^*$.
Then by using
\eqref{eq:Di_Gr} and
\eqref{eq:S_and_D_i}, we have
\begin{align*}
\iota\bigl(\Gro{w}(z|\eta)\bigr)&=
\iota\bigl(D_i^Q \Gro{s_iw}(z|\eta)\bigr)=D_{n-i}^Q \iota\bigl(\Gro{s_iw}(z|\eta)\bigr)\\
&=D_{n-i}^Q \Gro{(s_iw)^*}(z|\eta)=\Gro{s_{n-i}(s_iw)^*}(z|\eta)=\Gro{w^*}(z|\eta).\tag*{\qed}
\end{align*}\renewcommand{\qed}{}
\end{proof}

\subsection[Explicit formula for Gro\{s\_theta\}(x|eta)]{Explicit formula for $\boldsymbol{\Gro{s_\theta}(x|\eta)}$}
The quantum double Grothendieck polynomial for $s_\theta$ plays a~fundamental role
in the following.

Let $\omega^{(\eta)}$ be the cyclic permutation with respect to the variables $\eta$, that is
$\omega^{(\eta)}(\eta_i)=\eta_{i+1}$
with $\eta_{n+1}=\eta_1$.
We use the following notation:
$
[x|\eta]^i=(x\oplus \eta_1)\cdots (x\oplus \eta_i)
$
for $i\ge 1$.
\begin{Proposition}
\label{prop:Gtheta}
We have
\begin{equation}
\Gro{s_\theta}(z|\eta)=\Gro{s_1s_2\dots s_{n-1}}(z|\eta)\cdot \Gro{s_{n-2}\dots s_2s_1}\bigl(z|\omega^{(\eta)}(\eta)\bigr)\label{eq:Gtheta}
.
\end{equation}
\end{Proposition}
\begin{proof}
We first prove
the classical version
\begin{equation}
\mathfrak{G}_{s_\theta}(z|\eta)=\mathfrak{G}_{s_1s_2\dots s_{n-1}}(z|\eta)\cdot \mathfrak{G}_{s_{n-2}\dots s_2s_1}\bigl(x|\omega^{(\eta)}(\eta)\bigr)
\label{eq:Gtheta-cl}
\end{equation}
of
\eqref{eq:Gtheta}.
Since
$s_\theta\in S_n$ is a~dominant permutation of
code (see \cite[Chapter~I]{Mac} for the definitions) $c(s_\theta)=(n-1,1,\dots,1,0)$,
\begin{align*}
\mathfrak{G}_{s_\theta}(z|\eta)&=
[x_1|\eta]^{n-1}
[x_2|\eta]
\cdots
[x_{n-1}| \eta]=[x_1|\eta]
[x_2|\eta]
\cdots
[x_{n-1}| \eta]\times
\big[x_1|\omega^{(\eta)}(\eta)\big]^{n-2}.
\end{align*}
For the first equality, we have used \cite[Lemma B.6]{MNS1}.
Because
$s_1s_2\dots s_{n-1}$
and $s_{n-2}\dots s_{2}s_1$
are also dominant, with the codes
$(1,\dots,1,0)$ and $ (n-2,0,\dots,0)$
respectively, we have by \cite[Lemma B.6]{MNS1}
\begin{gather*}%\label{eq:phi_(n-1)}
\mathfrak{G}_{s_1s_2\dots s_{n-1}}(z|\eta)
=[x_1|\eta]
[x_2|\eta]
\cdots
[x_{n-1}| \eta],\qquad
\mathfrak{G}_{s_{n-2}\dots s_{2}s_1}(z|\eta)
=[x_1|\eta]^{n-2}.
\end{gather*}
Therefore, \eqref{eq:Gtheta-cl} holds.

We will apply the quantization map
$\hat{Q}$ of Lenart--Maeno (see \cite[Section~3]{LM}) to~\eqref{eq:Gtheta-cl} with respect to the variables
$x_1,\dots,x_n$, where $x_i=1-z_i$.
Let \smash{$f_j^{(i)}=e_j(1-x_1,\dots,1-x_i)$}. Then
$\hat{Q}$ is a~linear map such that
\[\hat{Q}\bigl(f_{p_1}^{(1)}f_{p_2}^{(2)}\cdots f_{p_{n-2}}^{(n-2)}f_{p_{n-1}}^{(n-1)}\bigr)=F_{p_1}^{(1)}F_{p_2}^{(2)}\cdots F_{p_{n-2}}^{(n-2)}F_{p_{n-1}}^{(n-1)}
\]
for $0\le p_i\le i$ (see \cite[Proposition 3.16]{LM}).
As a~polynomial in $x_1$, $[x_1|\eta]^{n-2}$
has degree~${n-2}$.
Hence it is easy to see that
$\mathfrak{G}_{s_{n-2}\dots s_{2}s_1}(z|\eta)=[x_1|\eta]^{n-2}$
is a~linear combination of
$\smash{f_{p_1}^{(1)}f_{p_2}^{(2)}}\cdots\allowbreak\smash{ f_{p_{n-2}}^{(n-2)}}$ with $ 0\le p_i\le i$.
Also,
\begin{equation}
\mathfrak{G}_{s_{1}s_2\dots s_{n-1}}(z|\eta)=[x_1|\eta]
[x_2|\eta]
\cdots
[x_{n-1}| \eta]=
1+\sum_{j=1}^{n-1}
(\eta_1-1)^jf_j^{(n-1)}.
\label{eq:phi_(n-1)-2}
\end{equation}
Therefore,
$\hat{Q}$ preserves the product of~\eqref{eq:Gtheta-cl},
so we have~\eqref{eq:Gtheta}.
\end{proof}

By the proof of the previous
proposition, we have the following.
\begin{Corollary}\label{cor:phi_(n-1)}
We have
\[
\Gro{s_1s_2\dots s_{n-1}}(z|\eta)=
\sum_{j=0}^{n-1}(-1)^j(1-\eta_{1})^{j}F_j^{(n-1)}.
\]
\end{Corollary}
\begin{proof} This identity is obtained by
applying $\hat{Q}$ to
\eqref{eq:phi_(n-1)-2}.
\end{proof}





\subsection[Affine K-nil-Hecke action on the equivariant quantum K-ring]{Affine $\boldsymbol{K}$-nil-Hecke action on the equivariant quantum $\boldsymbol{K}$-ring}
Here we explain an action of
the $0$-th Demazure operator on
$QK_T(\SL_n(\C)/B)_Q$.
We define an operator \smash{$D_0^Q$} on
$QK_T(\SL_n(\C)/B)_Q$ by
\begin{equation}
\label{eq:D0Q}
D_0^Q=T^{Q}_{\theta}+
Q^{-\theta^\vee}
\mathscr{O}^{s_\theta}\cdot
s_\theta,
\end{equation}
where an element $T^Q_\theta$ is defined by
\[
T^Q_\theta:=\bigl(\omega^{(\eta)}\bigr)^{-i}\circ T^Q_i\circ \bigl(\omega^{(\eta)}\bigr)^i=\frac{s_\theta^{(\eta)}-1}{\eta_1\ominus\eta_n}=\frac{s^{(\eta)}_\theta-1}{1-{\rm e}^{-\theta}}.
\]

If we present $QK_T(\SL_n(\C)/B)$ as a~quotient ring by~\cite{MNS1}, we can replace
$Q^{-\theta^\vee}\mathscr{O}^{s_\theta}$, acting as a~multiplication operator,
with \smash{$Q^{-\theta^\vee}\Gro{s_\theta}$}.

Together with $D_i^Q$, $i\in \Iaf$, and the left multiplication by $R(T)$,
 $QK_T(\SL_n(\C)/B)_Q$ has a~structure of $K$-nil-DAHA-module (see~\cite{KNS,KoNOS,Orr}).


\begin{Proposition}\label{prop:Quantum_D(1)}
Let $x=wt_\xi\in \hat{W}_G^0$. Then, with the presentation~{\rm\cite{MNS1}}, we have
\[%\label{eq:Quantum_D(1)}
D_x^Q(1)
=Q^\xi \Gro{w} \mod \mathscr{I}_n^Q.
\]
\end{Proposition}
See Appendix~\ref{sec:nil-Hecke_on_QK}.



\subsection{Proof of Theorem~\ref{thm:main}}
The purpose of this section is to prove Theorem~\ref{thm:main}
\begin{Remark}
The image of $\tilde{\Phi}_n$ is \smash{$\hat{\Lambda}_{(n)}^\Rep\big[\sigma_i^{-1},(\sigma(\sigma_i))^{-1}|1\leq i\leq n\big]$}.
\end{Remark}





The following result is an important special case of Theorem~\ref{thm:main}.
This proposition in turn implies the commutativity of $\Phi_n$ with the Demazure operators (see Corollary~\ref{cor:Phi_commutes_w_Dem}).

\begin{Proposition}
\label{prop:key} We have
\[
\tilde{\Phi}_n\bigl(Q^{-\theta^\vee}\Gro{s_\theta}\bigr)=\tilde{g}_{s_0}^{(k)}(y|b).
\]
\end{Proposition}
\begin{Remark}The proof of this result is purely combinatorial.
We do not have to resort to
the deep geometric fact of Corollary~\ref{cor:BF}.
\end{Remark}

We first compute
the image of
the factor
$\Gro{s_1s_2\dots s_{n-1}}$ of $\Gro{s_\theta}$
in the factorized form~\eqref{eq:Gtheta} under~$\Phi_n$.
\begin{Lemma}\label{lem:Phi_1st_factor}
We have
\begin{gather}
 (\Phi_n\circ \iota)\bigl(\Gro{s_1s_2\dots s_{n-1}}\bigr)=
\frac{{\rm e}^{-a_1}z_{11}}{\tau_1},\label{eq:Phi_1st_factor}
 \\
 \Phi_n\bigl(\Gro{s_1s_2\dots s_{n-1}}\bigr)=
 \frac{\tau_n}{\tau_{n-1}}
 \frac{{\rm e}^{a_n}}{z_{nn}}.\label{eq:Phi_1st_factor-2}
\end{gather}
\end{Lemma}
\begin{proof}
 We will show~\eqref{eq:Phi_1st_factor}. Then
 \eqref{eq:Phi_1st_factor-2} is obtained from this by Proposition~\ref{prop:S_on_tau}.
Since $\iota(1-\eta_{1})={\rm e}^{-a_1}$, we have by Corollary~\ref{cor:phi_(n-1)}
\[
\iota\bigl(\Gro{s_1s_2\dots s_{n-1}}\bigr)=
\sum_{j=0}^{n-1}(-1)^j{\rm e}^{-ja_1}\iota\bigl(F_j^{(n-1)}\bigr).
\]
 By Lemma~\ref{lem:SF_to_c}, we have
\smash{$(\Phi_n\circ \iota)
\bigl(F^{(n-1)}_j\bigr)=(\iota\circ \Phi_n)
\bigl(F^{(n-1)}_j\bigr)=(-1)^j{c_j^{(1)}}/{\tau_1}$}.
Note that we have~\smash{$\tau_1=c_0^{(1)}$} from~\eqref{eq:det_c_to_tausigma}.
Therefore, we have
\begin{align*}
(\Phi_n\circ \iota)
\bigl(\Gro{s_1s_2\dots s_{n-1}}\bigr)
&=\frac{1}{\tau_1}
\sum_{j=0}^{n-1}{\rm e}^{-ja_1}
c_j^{(1)}=\frac{{\rm e}^{-a_1}z_{11}}{\tau_1}
\qquad \text{by~\eqref{eq:zii_c_2}}.\tag*{\qed}
\end{align*}\renewcommand{\qed}{}
\end{proof}
 In order to compute the image of
 $\Gro{s_{n-2}\dots s_2s_{1}}\bigl(z|\omega^{(\eta)}\eta\bigr)$ under $\Phi_n$, we prepare
 the following lemma.




\begin{Proposition}\label{prop:Phi_n_2nd_factor}
We have
\[%\label{eq:Phi_n_2nd_factor}
\Phi_n\bigl(\Gro{s_{n-2}\dots s_2s_{1}}(z|\eta)\bigr)=\frac{{\rm e}^{-a_1-a_2}z_{12}}{\tau_1}.
\]
\end{Proposition}
\begin{proof}
By~\eqref{eq:Di_Gr} and Proposition~\ref{prop:S_star}, we have
\begin{equation}
\Gro{s_{n-2}\dots s_2s_{1}}(z|\eta)
=
\iota\bigl(D_1^{Q}\Gro{s_{1}s_2\dots s_{n-1}}(z|\eta)\bigr).\label{eq:2nd_factor}
\end{equation}
Applying
${\Phi}_n$ to both sides of~\eqref{eq:2nd_factor}, we have
\begin{align*}
\Phi_n\bigl(\Gro{s_{n-1}\dots s_2s_{1}}(z|\eta)\bigr)
&=(\Phi_n\circ \iota)\bigl(D_1^{Q}\Gro{s_{1}s_2\dots s_{n-1}}(z|\eta)\bigr) \\
&=
D_1\bigl(
(\Phi_n\circ \iota)\bigl(\Gro{s_{1}s_2\dots s_{n-1}}(z|\eta)\bigr)
\bigr)\qquad\text{by~\eqref{eq:S_and_D_i} and~\eqref{eq:D_and_Phi}}\\
&=D_1\left(
\frac{{\rm e}^{-a_1}z_{11}}{\tau_1}
\right)\qquad \text{by
\eqref{eq:Phi_1st_factor}}\\
&=
\frac{1}{\tau_1}
D_1({{\rm e}^{-a_1}z_{11}})
\qquad \text{since $\tau_1$ is $S_n$-invariant.}
\end{align*}
Finally, it is easy to show
$D_1({\rm e}^{-a_1}z_{11})={{\rm e}^{-a_1-a_2}z_{12}}$
by using~\eqref{eq:cent}.
\end{proof}

\begin{proof}[Proof of Proposition~\ref{prop:key}]
Note that $\omega^{(\eta)}$ corresponds to $\omega^{-1}$
under the identification~\eqref{eq:eta-b}.
By using Proposition~\ref{prop:Gtheta}, \eqref{eq:Phi_1st_factor-2}, and
Proposition~\ref{prop:Phi_n_2nd_factor}, we have
\begin{align*}
\Phi_n
\bigl(\Gro{s_{\theta}}(z|\eta)\bigr)
&=
\Phi_n
\bigl(\Gro{s_1s_2\dots s_{n-1}}(z|\eta)\cdot \Gro{s_{n-2}\dots s_2s_1}\bigl(x|\omega^{(\eta)}(\eta)\bigr)
\bigr)\\
&=
\frac{\tau_n}{\tau_{n-1}}
\frac{{\rm e}^{a_n}}{z_{nn}}
\cdot
\omega^{-1}
\left(
\frac{{\rm e}^{-a_1-a_2}z_{12}}{\tau_1}
\right)\\
&=
\frac{\tau_n}{\tau_{n-1}}
\frac{{\rm e}^{a_n}}{z_{nn}}
\cdot
\frac{{\rm e}^{-a_n-a_1}z_{n,n+1}}{\tau_1}
\qquad\text{
since $\omega(z_{ij})=z_{i+1,j+1}$}
\\
&=
\frac{\tau_n}{\tau_{n-1}\tau_1}
\frac{{\rm e}^{-a_1}z_{n,n+1}}{z_{nn}}.
\end{align*}
From~\eqref{eq:zij_in_KkSchur},
we have
$z_{n,n+1}=\omega^{-1}(z_{12})={\rm e}^{a_n}z_{nn}g_{s_0}^{(k)}(y|b)$.
Thus, together with~\eqref{eq:sigma(tg)}, we have
\[
\frac{{\rm e}^{-a_1}z_{n,n+1}}{z_{nn}}
={\rm e}^{-a_1+a_n}g_{s_0}^{(k)}(y|b)
=\sigma^{-1}\bigl(\tilde{g}_{s_0}^{(k)}(y|b)\bigr).
\]
In view of
$
\tilde{\Phi}_n\bigl(Q^{\theta^\vee}\bigr)=
\frac{\tau_n}{\tau_{n-1}\tau_1}$,
the proof is complete.
\end{proof}







Now we can prove a~crucial property of $\tilde{\Phi}_n$.
\begin{Proposition}\label{prop:Phi_D0} We have
$\tilde{\Phi}_n\circ D_0^{Q}=D_0\circ\tilde{\Phi}_n
$.
\end{Proposition}
\begin{proof}
Since $s_\theta^*=s_\theta$, we have $\tilde{\Phi}_n\circ T_\theta^Q
=T_\theta \circ \tilde{\Phi}_n$
and $\tilde{\Phi}_n\circ s_\theta
=s_\theta \circ \tilde{\Phi}_n$.
From Proposition~\ref{prop:key}, we have
\begin{align*}
\tilde{\Phi}_n\circ D_0^{Q}&=\tilde{\Phi}_n
\circ \bigl(T_\theta^{Q}+
Q^{-\theta^\vee}
\Gro{s_\theta}
s_\theta
\bigr)=\tilde{\Phi}_n\circ T_\theta^{Q}+
\tilde{\Phi}_n\circ\bigl(Q^{-\theta^\vee}
\Gro{s_\theta}(x,y)
s_\theta\bigr)\\
&=T_\theta\circ\tilde{\Phi}_n
+
\tilde{\Phi}_n\bigl(Q^{-\theta^\vee}
\Gro{s_\theta}(x,y)\bigr)
\tilde{\Phi}_n\circ
s_\theta=T_\theta\circ\tilde{\Phi}_n+
\tilde{g}_{s_0}^{(k)}(y|b)\cdot
\tilde{\Phi}_n\circ
s_\theta
\\
&=T_\theta\circ\tilde{\Phi}_n+
D_0(1)\cdot
s_\theta\circ \tilde{\Phi}_n=D_0\circ\tilde{\Phi}_n,
\end{align*}
where the last equality follows from the simply identity
$
D_0=T_\theta
+ D_0(1)
s_\theta$.
\end{proof}


\begin{Corollary}\label{cor:Phi_commutes_w_Dem}
Let \smash{$x\in \hat{W}_G^0$}. Then
$
\tilde{\Phi}_n\circ D_x^Q=
D_{x^{\omega_k}}
\circ \tilde{\Phi}_n$.
\end{Corollary}
\begin{proof}
This is an immediate consequence of Proposition~\ref{prop:Phi_D0} and~\eqref{eq:D_and_Phi}.
\end{proof}



\begin{proof}[Proof of Theorem~\ref{thm:main}]
 \begin{align*}
 \tilde{\Phi}_n\bigl(Q^\xi \Gro{w}\bigr)
&=
\tilde{\Phi}_n\bigl(D_x^Q(1)\bigr) \qquad \text{by Proposition~\ref{prop:Quantum_D(1)}}\\
&={D}_{x^{\omega_k}}\bigl(\tilde{\Phi}_n(1)\bigr) \qquad \text{by Corollary~\ref{cor:Phi_commutes_w_Dem}}\\
&={D}_{x^{\omega_k}}(1)=\tilde{g}_{x^{\omega_k}}^{(k)}(y|b).\tag*{\qed}
\end{align*} \renewcommand{\qed}{}
\end{proof}


\section[Determinantal formulas for the k-small g\_la\^{}(k)(y|b)]{Determinantal formulas for the $\boldsymbol{k}$-small $\boldsymbol{\tilde{g}_\la^{(k)}(y|b)}$}\label{sec:det}

\subsection[S\_n-and T\_i-actions on O(Z)]{$\boldsymbol{S_n}$-and $\boldsymbol{T_i}$-actions on $\boldsymbol{\mathscr{O}(\mathscr{Z})}$}

Let $s_i$ $(1\le i\le n-1)$ be the generators of the symmetric group $S_n$, which acts on $\mathscr{O}\bigl(\tilde{\mathscr{Z}}\bigr)^{\reg}$ as a~$\CC$-algebra morphism by
\[
s_i\bigl({\rm e}^{\pm a_j}\bigr)={\rm e}^{\pm a_{s_i(j)}},\qquad
s_i(Z)=\mathfrak{s}_iZ\mathfrak{s}_i^{-1},
\qquad \mbox{where} \quad \mathfrak{s}_i=E-\bigl({\rm e}^{-a_i}-{\rm e}^{-a_{i+1}}\bigr)E_{i+1,i}.
\]

The $S_n$-action preserves the subalgebra $\mathscr{O}({\mathscr{Z}})$.
The isomorphism \smash{$\beta\colon\mathscr{O}(\mathscr{Z})\to \hat{\Lambda}_{(n)}^\Rep$} given in Theorem~\ref{thm:ISY} satisfies $s_i\circ\beta=\beta\circ s_i$ for all $i$ because $s_i(Z)$ is contained in the centralizer of~${s_i(A):=A|_{a_i\leftrightarrow a_{i+1}}}$.

On the coordinate functions $z_{ij}$, $S_n$ acts as
\begin{gather}
s_i(z_{jj})=z_{s_i(j),s_i(j)}\label{eq:s-2},\\
s_i(z_{ki})=z_{ki}+\bigl({\rm e}^{-a_i}-{\rm e}^{-a_{i+1}}\bigr)z_{k,i+1},\qquad k<i,\label{eq:s-3}\\
s_i(z_{i+1,j})=z_{i+1,j}-\bigl({\rm e}^{-a_i}-{\rm e}^{-a_{i+1}}\bigr)z_{ij},\qquad j>i,\label{eq:s-4}\\
s_i(z_{kj})=z_{kj},\qquad
k\ne i+1 \quad\text{or}\quad j\ne i.
\label{eq:s-5}
\end{gather}
It should be noted that~\eqref{eq:s-2} is a~consequence of~\eqref{eq:s-3}, and $z_{ij}$ is defined by~\eqref{eq:cent}.

\begin{Proposition}\label{prop:S_n_invariant_functions}
The action variables $c_0,c_1,\dots,c_{n-1}$ are $S_n$-invariant.
\end{Proposition}
\begin{proof}
From~\eqref{eq:def_of_c_i} and the definition of $s_i$, we have
\begin{align*}
s_i(Z)&=\mathfrak{s}_iZ\mathfrak{s}_i^{-1}
=\mathfrak{s}_i\left(\sum_{j=0}^{n-1}c_jA^j\right)\mathfrak{s}_i^{-1}\\
&=\sum_{j=0}^{n-1}c_j\mathfrak{s}_iA^j\mathfrak{s}_i^{-1}=\sum_{j=0}^{n-1}c_j\bigl(\mathfrak{s}_iA\mathfrak{s}_i^{-1}\bigr)^j=\sum_{j=0}^{n-1}c_js_i(A)^j.
\end{align*}
Since $s_i(A)$ is also regular, this equation uniquely determines $c_j$'s.
On the other hand, because~$s_i$ is a~$\C$-algebra homomorphism we have
$s_i(Z)=\sum_{j=0}^{n-1}s_i(c_j)s_i(A)^j$.
Hence we have $s_i(c_j)=c_j$.
\end{proof}
Theorem~\ref{thm:det_c_to_tausigma} implies the following important consequence.
\begin{Corollary}\label{cor:tausigma invariant}
 $\tau_i$ and $\sigma_i$ are $S_n$-invariant.
\end{Corollary}
\begin{Remark}\label{rem:PinvZP}
We have
\[
P^{-1}ZP=\sum_{j=0}^{n-1}c_j C_A^j.
\]
Since the entries of $C_A$ are $S_n$-invariant,
the entries of $P^{-1}ZP$ are $S_n$-invariant by Proposition~\ref{prop:S_n_invariant_functions}.
\end{Remark}

\begin{Proposition}\label{prop:T_action_on_z}
For $1\le i\le n-1$,
\begin{gather}
T_i(z_{ki})={\rm e}^{-a_i}z_{k,i+1},\label{eq:z_Tz}\\
T_i(z_{i+1,j})=-{\rm e}^{-a_i}z_{ij},\label{eq:z_Tz_2}\\
T_i(z_{kj})=0,\qquad
k\ne i+1 \quad \text{or}\quad j\ne i.\label{eq:z_Tz_3}
\end{gather}
\end{Proposition}

\begin{proof}
Equations~\eqref{eq:z_Tz}, \eqref{eq:z_Tz_2}, and~\eqref{eq:z_Tz_3} are immediate consequences of~\eqref{eq:s-3}, \eqref{eq:s-4}, and~\eqref{eq:s-5}, respectively.
\end{proof}

For convenience, we extend the definition of $z_{ij}$ to any $1\le i\le j$ by letting $z_{ii}=z_{i+n,i+n}=z_{i+2n,i+2n}=\cdots$ and
\[%\label{eq:def_of_z_extended}
z_{ij}=-
\frac{z_{i,j-1}-z_{i+1,j}}{{\rm e}^{-a_{i(\mathrm{mod}\, n)}}-{\rm e}^{-a_{j(\mathrm{mod}\, n)}}}.
\]
By definition, $z_{ij}$ is an element of $\mathscr{O}\bigl(\Tilde{\mathscr{Z}}\bigr)^{\reg}$.
Let $\omega:=s_1s_2\dots s_{n-1} $.
It is shown by induction on~${j-i\geq 0}$ that
\begin{equation}\label{eq:pi_zij}
\omega(z_{ij})=z_{i+1,j+1},
\end{equation}
which implies $z_{ij}\in \mathscr{O}(\mathscr{Z})$.
In particular, we have $z_{ij}=z_{i+n,j+n}$.

\begin{Remark}
$T_\theta$ can be seen as the $n$-th divided difference operator.
Proposition~\ref{prop:T_action_on_z} is naturally extended to
the case for $i=n$ as
$
T_\theta(z_{k,n})={\rm e}^{-a_n}z_{k,n+1}$ and $ T_\theta(z_{n+1,j})=-{\rm e}^{-a_n}z_{nj}$.
Moreover, the expressions~\eqref{eq:s-2}--\eqref{eq:z_Tz_2} are also valid for arbitrarily $i,k,j\in \ZZ$ under the identification $T_{i+n}=T_i$ and $T_n=T_\theta$.
\end{Remark}

\subsection[Determinantal formula for g\^{}(k)\_lambda]{Determinantal formula for $\boldsymbol{\tilde{g}^{(k)}_\lambda}$}
\label{sec:det_formula}
A $k$-bounded partition $\lambda$ is \textit{$k$-small} if $\lambda$ is contained in at least one of $R_1,R_2,\dots,R_{n-1}$.
This is equivalent to $\ell(\la)+\la_1\le n$.

\subsubsection{Notation}
We identify the set $\Iaf=\{0,1,\dots,n-1\}$ of
the type \smash{$A_{n-1}^{(1)}$} affine Dynkin nodes
with $\Z/n\Z$.
The $n$-\emph{residue} is the map
$
\mathrm{res}\colon
\mathbb{N}\times \mathbb{N}
\rightarrow \Iaf\cong \Z/n\Z$, $
(i,j)\to j-i \mod n$.
For $x=(i,j)\in \la$, let~${
\mathfrak{d}(x):=\{{\rm e}^{-a_{\mathrm{res}(s,j)}}\mid i\le s\le \la_j'\}}$.


For any subset $X$ of $\{{\rm e}^{-a_1},\dots,{\rm e}^{-a_n}\}$
we denote by $h_m(X)$ the $m$-th complete symmetric
polynomial in $X$.
We use
the abbreviation
\smash{$
\ha_m^{\langle i,j\rangle;\la}:=h_m(\mathfrak{d}(i,j))$},
for $(i,j)\in \la$.
For $k$-small partition $\la$, define
\begin{equation}
\label{eq:def xi}
\co_\la(y)=\prod_{i=1}^{\la_1'}\Omega(b_{\mathrm{res}(i,1)}|y).
\end{equation}




\begin{Example}\label{ex:h d}
When $n=6$ and
$
\lambda=(3,3,1)$.
If we fill the boxes of $\la$ with the $n$-residue, we have
\[
\begin{ytableau}
0&1&2\\
5&0&1\\4
\end{ytableau}.
\]
For example, $\mathfrak{d}(2,1)=\big\{{\rm e}^{-a_5},{\rm e}^{-a_4}\big\}$ so
\smash{$\ha_m^{\langle 2,1\rangle;\la}=h_m({\rm e}^{-a_5},{\rm e}^{-a_4})$}.
We also have
$
\co_\la(y)=\Omega(b_6|y)\allowbreak\times\Omega(b_5|y)\Omega(b_4|y)$.
\end{Example}



\subsubsection[Determinantal formula for g\_la(y|b) for k-small la]{Determinantal formula for $\boldsymbol{\tilde{g}_\la(y|b)}$ for $\boldsymbol{k}$-small $\boldsymbol{\la}$}


\begin{Theorem}\label{thm:determinant_formula_for_ksmall_g}
Let $\lambda$ be a~$k$-small partition.
Set $
l=n-\la_1'+1$
and $r=\la_1$.
Define the following
square matrix of size $n-l+1+r$
\begin{equation}\label{eq:M la}
 M_\la=
 \left(
\begin{array}{@{}cccc|cccc@{}}
 z_{ll}&\cdots &\cdots &z_{ln}&z_{l,n+1}&\cdots&\cdots&z_{l,n+r}\\
 &\ddots& &\vdots&\vdots&\cdots&\cdots&\vdots\\
 &&\ddots &\vdots&\vdots&\cdots&\cdots&\vdots\\
 & &&z_{nn}&z_{n,n+1}&\cdots&\cdots&z_{n,n+r}\\ \hline
 &\cdots
 %-h_3^{\bl_{31}}
 & \ha_2^{\langle {2,1}\rangle;\la} & -\ha_1^{\langle 1,1\rangle;\la} & 1 & & \\
 &&\cdots
 %-\ha3^{\bl(3,2)}
 &\ha_2^{\langle 2,2\rangle;\la}&-\ha_1^{\langle 1,2\rangle;\la} & \ddots &&\\
 &&\cdots&\vdots&\vdots & \ddots &\ddots&\\
 &&&\pm \ha_r^{\langle r,r\rangle;\la}&\mp \ha_{r-1}^{\langle r-1,r\rangle;\la}&\cdots &-\ha_1^{\langle 1,r\rangle;\la} &1
\end{array}
\right).
\end{equation}
Then we have
\begin{equation}\label{eq:matrix_M}
\tilde{g}^{(k)}_{\lambda}(y|b)=
\frac{\det(M_\lambda)}{\co_\la(y)}.
\end{equation}
\end{Theorem}


\begin{Remark}\quad
\begin{enumerate}
\itemsep=0pt
\item[(1)] $l$ is the $n$-residue of the bottom box of the first column of $\la$.
 \item[(2)] We have $\la_1<l$ because $\la$ is $k$-small.
If we take $r$ such that $\la_1\le r<l$ and
 consider the matrix $M_\la$ by the same formula~\eqref{eq:M la},
 then~\eqref{eq:matrix_M} also holds.
\end{enumerate}

\end{Remark}



\begin{Example}\label{ex:k small det}
Let $n=6$ and $\la$ be as in Example~\ref{ex:h d}. We have $l=4$.
The formula for $\tilde{g}^{(5)}_{\lambda}(y|b)$ by taking $r=3$ is
\[
\frac{1}{\Omega(b_4|y)\Omega(b_5|y)\Omega(b_6|y)}
\begin{vmatrix}
 z_{44}&z_{45}&z_{46}&z_{47}&z_{48}&z_{49}\\
 0&z_{55}&z_{56}&z_{57}&z_{58}&z_{59}\\
 0&0&z_{66}&z_{67}&z_{68}&z_{69}\\
 -f_3^{\langle 3,1\rangle;\la} & f_2^{\langle 2,1\rangle;\la} & -f_1^{\langle 1,1 \rangle;\la} & 1 & 0 & 0\\
 0&0&f_2^{\langle 2,2\rangle;\la}&-f_1^{\langle 1,2 \rangle;\la} & 1 &0\\
 0&0&0&f_2^{\langle 2,3\rangle;\la}&-f_1^{\langle 1,3 \rangle;\la} & 1
\end{vmatrix}.
\]
\end{Example}

\begin{Remark}
\bgroup
\renewcommand\binom[2]{\genfrac{[}{]}{0pt}{1}{#1}{#2}}
In the non-equivariant case $a_i=0$, \eqref{eq:matrix_M} reduces to
\[
\Tilde{g}^{(k)}_\lambda(y|0)=%\det
\begin{vmatrix}
h_0(y)&h_1(y)&h_2(y)&\cdots&\cdots&\cdots &\cdots &h_{n-1}(y)\\
&h_0(y)&h_1(y)&h_2(y)&\cdots&\cdots &\cdots &h_{n-2}(y)\\
&&\ddots&\ddots&\ddots&\cdots &\cdots &\vdots\\
&&&h_0(y)&h_1(y)&h_2(y)&\dots &h_r(y)\\
\binom{\lambda_1'}{\lambda_1'}&\cdots&\binom{\lambda_1'}{2}&\binom{\lambda_1'}{1}&\binom{\lambda_1'}{0}\\
&\binom{\lambda_2'}{\lambda_2'}&\cdots&\binom{\lambda_2'}{2}&\binom{\lambda'_2}{1}&\binom{\lambda_2'}{0}\\
&&\ddots&\cdots&\cdots&\ddots&\ddots\\
&&&&\binom{\lambda_r'}{\lambda_r'}&\cdots&\binom{\lambda_r'}{1}&\binom{\lambda_r'}{0}
\end{vmatrix},
\]
where $h_i=h_i(y)$ and $\binom{a}{b}=(-1)^b\genfrac{(}{)}{0pt}{1}{a}{b}$.
\egroup
\end{Remark}

The proof of Theorem~\ref{thm:determinant_formula_for_ksmall_g} is divided into two cases: (i) the case when $\lambda$ is a~one-column partition and (ii) the case when $\lambda$ is a~general $k$-small partition.

\subsubsection[Case: la=(1\^{}i)]{Case: $\boldsymbol{\la=\bigl(1^i\bigr)}$}
Let $\lambda$ is a~one-column partition $\lambda=\bigl(1^i\bigr)$.
We have
\[
 M_{(1^i)}=
 \left(
\begin{array}{@{}cccc|c@{}}
 z_{n-i+1,n-i+1}&\cdots &\cdots &z_{n-i+1,n}&z_{n-i+1,n+1}\\
 &\ddots& &\vdots&\vdots\\
 &&\ddots &\vdots&\vdots\\
 & &&z_{nn}&z_{n,n+1}\\ \hline
 (-1)^i\ha_i^{\langle {i,1}\rangle;(1^i)} &\cdots
 %-3^{\bl_{31}}
 & \ha_2^{\langle {2,1}\rangle;(1^i)} & -\ha_1^{\langle 1,1\rangle;(1^i)}&1
\end{array}
\right).
\]

In this case, \smash{$\det(M_{(1^i)})$} given in~\eqref{eq:matrix_M} is expanded as
\[%\label{eq:expand_M}
\det(M_{(1^i)})
=
\sum_{m=0}^{i}
\ha_{m}^{\langle m,1\rangle;(1^i)}
d^{(i)}_m \\
\]
by the expansion along the bottom row,
where
\smash{$%\label{E:dm}
d^{(i)}_m=\det Z_{[n-i+1,n]}
^{[n-i+1,n+1]\setminus \{n-m+1\}}$}.
For $0\le m\le i< n$, define
\begin{gather}
f^{(i)}_m=
\ha_{m}^{\langle m,1\rangle; (1^i)}
\cdot d^{(i)}_m.\label{eq: varphi m}
\end{gather}
Then~\eqref{eq:matrix_M} reads
\begin{equation}\label{eq:$K$-$k$-one_column_determinant}
\begin{aligned}
\tilde{g}^{(k)}_{(1^i)}(y|b)
=
\frac{
\sum_{m=0}^{i}
\varphi_m^{(i)}
}{\co_{(1^i)}(y)}
\end{aligned}.
\end{equation}

We consider the minor
$d^{(i+1)}_m$
of size $i+1$.
Note that we have
\begin{equation}
 d^{(i+1)}_m=z_{n-i,n-i}d^{(i)}_{m} \label{eq: d=dz}
\end{equation}
since $Z$ is upper triangular.
It follows that
\begin{equation}\label{eq:tilde varphi}
z_{n-i,n-i}\varphi_m^{(i)}=\ha_{m}^{\langle m,1\rangle; (1^i)}d^{(i+1)}_m.
\end{equation}


\begin{Lemma}\label{lem:T_on_f}
For $1\le i\le n-1$, we have \smash{$T_{n-i}\bigl(d^{(i+1)}_i\bigr)={\rm e}^{-a_{n-i}}d^{(i+1)}_{i+1}$}.
\end{Lemma}
\begin{proof}
From~\eqref{eq:z_Tz_2} and~\eqref{eq:z_Tz_3}, we have
\begin{equation}
\label{eq:T d i}
T_{n-i}\bigl(d^{(i)}_{i}\bigr)
=
T_{n-i}\left(\det Z_{[n-i+1,n]}^{[n-i+2,n+1]}\right)=-{\rm e}^{-a_{n-i}}\cdot \det Z_{[n-i,n]\setminus \{n-i+1\}}^{[n-i+2,n+1]}
\end{equation}
because every entry of the rows of \smash{$Z_{[n-i+1,n]}^{[n-i+2,n+1]}$} except for the first row is invariant under $s_{n-i}$.
Then by using this, we have
\begin{align*}
T_{n-i}\bigl(d^{(i+1)}_{i}\bigr)={}&T_{n-i}\bigl(z_{n-i,n-i}d^{(i)}_{i}\bigr)\qquad\text{by~\eqref{eq: d=dz}}\\
={}&T_{n-i}(z_{n-i,n-i})\cdot d^{(i)}_{i}
+s_{n-i}(z_{n-i,n-i})\cdot T_{n-i}\bigl(d^{(i)}_{i}\bigr) \\
={}&
{\rm e}^{-a_{n-i}}z_{n-i,n-i+1}\cdot
d^{(i)}_i
\\
&-z_{n-i+1,n-i+1}\cdot {\rm e}^{-a_{n-i}}
\det Z_{[n-i,n]\setminus \{n-i+1\}}^{[n-i+2,n+1]}\quad\text{by~\eqref{eq:z_Tz} and~\eqref{eq:T d i}}\\
={}&{\rm e}^{-a_{n-i}}d^{(i+1)}_{i+1},
\end{align*}
where the last equality follows by the expansion along the $1$st column of $d^{(i+1)}_{i+1}$.
\end{proof}
\begin{Lemma}\label{lem: D z varphi}
For $1\le i\le n-1$,
and $0\le m\le i$, we have
\[
D_{n-i}\bigl(z_{n-i,n-i}\varphi^{(i)}_m\bigr)=
\begin{cases}
\varphi^{(i+1)}_{m}, & m<i,\\
\varphi^{(i+1)}_{i}+\varphi^{(i+1)}_{i+1}, & m=i.
\end{cases}
\]
\end{Lemma}
\begin{proof}

Suppose $m<i$. We claim that $d^{(i+1)}_m$ is invariant under $s_{n-i}$.
In fact, we have
\begin{equation*}
 d^{(i+1)}_m=z_{n-i,n-i}z_{n-i+1,n-i+1}d^{(i-1)}_m
\end{equation*} by~\eqref{eq: d=dz},
which is $s_{n-1}$-invariant
by~\eqref{eq:s-2} and~\eqref{eq:z_Tz_3}.
We have
\begin{align*}
D_{n-i}\bigl({z_{n-i,n-i}{\varphi}}^{(i)}_m\bigr)&=
D_{n-i}\bigl(\ha_{m}^{\langle m,1\rangle; \bigl(1^i\bigr)}\cdot d^{(i+1)}_m\bigr)\qquad\text{by $\eqref{eq:tilde varphi}$}\\
&=D_{n-i}\bigl(\ha_{m}^{\langle m,1\rangle; \bigl(1^i\bigr)}\bigr)\cdot d^{(i+1)}_m\qquad\text{since $d^{(i+1)}_m$ is $s_{n-i}$-invariant}
\\&=
\ha_{m}^{\langle m,1\rangle; (1^{i+1})}\cdot d^{(i+1)}_m
\qquad\text{by Lemma~\ref{lem:T_on_h} }
\\&=\varphi^{(i+1)}_m.
\end{align*}

When $m=i$, noting that
$\ha_i^{\langle i,1\rangle; (1^{i})}
={\rm e}^{-ia_{n-i+1}}$,
we have
\begin{align*}
D_{n-i}\bigl({z_{n-i,n-i}{\varphi}}^{(i)}_i\bigr)
={}&
D_{n-i}\bigl(
\ha_i^{\langle i,1\rangle; (1^{i})}
%{\rm e}^{-ia_{n-i+1}}
d^{(i+1)}_i\bigr)\qquad\text{by~\eqref{eq:tilde varphi}}\\
={}&
D_{n-i}\bigl(\ha_i^{\langle i,1\rangle; (1^{i})}\bigr)\cdot d^{(i+1)}_i
+{\rm e}^{-ia_{n-i}}\cdot T_{n-i}\bigl(d^{(i+1)}_i\bigr)\\
={}&
h_i({\rm e}^{-a_{n-i}},{\rm e}^{-a_{n-i+1}})\cdot d_{i}^{(i+1)}\\
&+{\rm e}^{-(i+1)a_{n-i}}d^{(i+1)}_{i+1}
\qquad\text{by Lemmas~\ref{lem:T_on_h} and~\ref{lem:T_on_f}}\\
={}&\ha_i^{\langle i,1\rangle; (1^{i+1})}\cdot d_{i}^{(i+1)}+
\ha_{i+1}^{\langle i+1,1\rangle;(1^{i+1})}\cdot d^{(i+1)}_{i+1}
\\
={}&\varphi^{(i+1)}_i+\varphi^{(i+1)}_{i+1} \qquad\text{by~\eqref{eq: varphi m}}.\tag*{\qed}
\end{align*} \renewcommand{\qed}{}
\end{proof}

\begin{proof}[Proof of Theorem~\ref{thm:determinant_formula_for_ksmall_g} for $\boldsymbol{\la=(1)^i}$]
We show~\eqref{eq:$K$-$k$-one_column_determinant} by induction on $i\geq 1$.
When $i=1$, the desired equation is shown directly as follows:
\begin{align*}
\tilde{g}^{(k)}_{(1)}(y|b)
&=D_0(1)=
1+\Omega(b_1|y) T_\theta\left(\frac{1}{\Omega(b_1|y)}\right)
=
1+\Omega(b_1|y) T_\theta\left(\frac{z_{nn}}{\Omega(b_1|y)\Omega(b_n|y)}\right)\\
&=1+\frac{T_\theta(z_{nn})}{\Omega(b_n|y)}
=\frac{z_{nn}+{\rm e}^{-a_n}z_{n,n+1}}{\Omega(b_n|y)}
=\frac{\varphi^{(1)}_0+\varphi^{(1)}_1}{\Omega(b_n|y)}
=\frac{\varphi^{(1)}_0+\varphi^{(1)}_1}{\co_{(1)}(y)}.
\end{align*}
Suppose~\eqref{eq:$K$-$k$-one_column_determinant} holds for some $1\le i\le n-2$.
By induction hypothesis,
\begin{align*}
\tilde{g}^{(k)}_{(1^{i+1})}(y|b)
&=D_{n-i}\left(
\frac{\sum_{m=0}^{i}\varphi^{(i)}_m}{\co_{(1^i)}(y)}
\right)=
D_{n-i}\left(
\frac{z_{n-i,n-i}\sum_{m=0}^{i}{{\varphi}}^{(i)}_m}{\Omega(b_{n-i}|y)\co_{(1^i)}(y)} \right)
\\
&=
\frac{
\sum_{m=0}^{i}D_{n-i}\left({z_{n-i,n-i}{\varphi}}^{(i)}_m\right) }{\Omega(b_{n-i}|y)\co_{(1^i)}(y)}=
\frac{
\sum_{m=0}^{i+1}\varphi^{(i+1)}_m }{\co_{(1^{i+1})}(y)}\qquad\text{by Lemma~\ref{lem: D z varphi}},
\end{align*}
where in the third equality we use
fact that $\Omega(b_{n-i}|y)\co_{(1^i)}(y)$ is $s_{n-i}$ invariant.
\end{proof}

\subsubsection[Case: general k-small la]{Case: general $\boldsymbol{k}$-small $\boldsymbol{\la}$}
Before proving Theorem~\ref{thm:determinant_formula_for_ksmall_g} for general $k$-small $\lambda$, we list a~few important properties of ${M}_\lambda$.




\begin{Lemma}\label{lem:comparing_determinants}
Let $\lambda$ be a~$k$-small partition.
Assume $\la$ is $i$-addable.
Let $\kappa$ be a~partition
obtained from $\lambda$ by adding a~box of $n$-residue $i$. We also assume $\kappa$ is $k$-small.
Then
$
 D_i(\det M_\lambda)=\det M_\kappa $ if $i\neq 0$, and $
D_\theta(\det M_\lambda)=\det M_\kappa$.
\end{Lemma}

\begin{proof}
We may assume $\la\ne \varnothing$.
The integers $l$, $r$ are defined in Theorem~\ref{thm:determinant_formula_for_ksmall_g}.
Let $\hat{x}:=x-l+1$ and $x^\dagger:=x-l+1+n$ for $x\in \ZZ$.
In $M_\lambda$, $z_{q,p}$ is at the $(\hat{q},\hat{p})$-th position (if exists).
For~${1\le s\le r}$, the
$\bigl(s^\dagger,\hat{p}\bigr)$-th entry of $M_\lambda$ is either a~constant or a~complete symmetric polynomial in ${\rm e}^{-a_{b}},{\rm e}^{-a_{b+1}},\dots,{\rm e}^{-a_{p}}$ where $b$ is the $n$-residue of the bottom box of the $s$-th column.

Let \smash{$M_\lambda^{(i)}$} be the matrix obtained from $M_\lambda$ by applying the following column or row operations:
\begin{itemize}\itemsep=0pt
 \item[(a)] Add ${\rm e}^{-a_{i}}$ times the $\big(\hat{i}+1\big)$-th column to the $\hat{i}$-th column.
 \item[(b)] If $l\le i<n$, subtract ${\rm e}^{-a_{i}}$ times the $\hat{i}$-th row from the $\big(\hat{i}+1\big)$-th row.
\end{itemize}
Since $\lambda$ is $k$-small and $\la$ is $i$-addable, there is
a unique integer $j$ such that the $n$-residue of the bottom box of the $j$-th column is $i + 1$.
In \smash{$M^{(i)}_\lambda$}, all the entries in rows except for the $j^\dagger$-th one are $s_{i}$-invariant because:
\begin{itemize}
\itemsep=0pt
 \item the $(\hat{q},\hat{p})$-th entry is
 \begin{gather*}
 -{\rm e}^{a_i-a_{i+1}}z_{i,i+1}\qquad \text{if} \ q=i+1,\quad p=i,\\
 z_{qi}+{\rm e}^{-a_i}z_{q,i+1}\qquad \text{if} \ q\neq i+1,\quad p=i,\\
 z_{i+1,p}-{\rm e}^{-a_i}z_{ip} \qquad \text{if} \ q=i+1,\quad p\neq i,\\
 z_{qp} \qquad \text{if} \ q\neq i+1,\quad p\neq i,
 \end{gather*}
One can verify that these entries are $s_i$-invariant from~\eqref{eq:s-2}--\eqref{eq:s-5}.
\item the $\bigl(s^\dagger,\hat{i}\bigr)$-th entry $(s\neq j)$ is either a~constant or a~polynomial of the form
\[
h_t({\rm e}^{-a_{b}},\dots,{\rm e}^{-a_i})-{\rm e}^{-a_i}h_{t-1}({\rm e}^{-a_{b}},\dots,{\rm e}^{-a_i},{\rm e}^{-a_{i+1}}),
\]
up to a~sign,
where $b$ is the $n$-residue of the bottom box of the $s$-th column, and $t>0$.
It is not difficult to show the polynomial above is
$s_i$-invariant.

\item the $\bigl(s^\dagger,\hat{p}\bigr)$-th entry $(s\neq j,\ p\neq i)$ is either a~constant or a~symmetric polynomial in~\smash{${\rm e}^{-a_{b}},{\rm e}^{-a_{b+1}},\dots,{\rm e}^{-a_p}$} with $b$ the $n$-residue of the bottom box of
the $s$-th column. Note that we have $b\ne i+1$.
Clearly, these are $s_i$-invariant.
\end{itemize}
Therefore, it suffices to compute the $D_i$-actions on the $j^\dagger$-th row of \smash{$M^{(i)}_\lambda$} only.
Set
$m=\lambda'_j$.
Then
$m=j-i+n-1$,
and the \smash{$\bigl(j^\dagger,\hat{p}\bigr)$}-th entry of \smash{$M^{(i)}_\lambda$} is
\begin{gather*}
(-1)^{m-p+i+1}h_{m-p+i+1}\bigl({\rm e}^{-a_{i+1}},\dots,{\rm e}^{-a_p}\bigr) \qquad \text{if} \ \hat{p}<\hat{i},\\
(-1)^{m+1}{\rm e}^{-a_i}h_{m}\bigl({\rm e}^{-a_{i+1}}\bigr) \qquad \text{if} \ \hat{p}=\hat{i},\\
0 \qquad \text{if} \ \hat{p}>\hat{i}
\end{gather*}
and the $\bigl(j^\dagger,\hat{p}\bigr)$-th entry of $M^{(i)}_\kappa$ is
\begin{gather*}
(-1)^{m-p+i+1}h_{m-p+i+1}({\rm e}^{-a_i},{\rm e}^{-a_{i+1}},\dots,{\rm e}^{-a_p}) \qquad \text{if} \
\hat{p}<\hat{i},\\
(-1)^{m+1}
(
h_{m+1}({\rm e}^{-a_i})
-{\rm e}^{-a_i}h_{m}({\rm e}^{-a_i},{\rm e}^{-a_{i+1}})
) \qquad \text{if} \  \hat{p}=\hat{i},\\
0   \qquad \text{if} \ \hat{p}>\hat{i}.
\end{gather*}
By Lemma~\ref{lem:T_on_h}, the image of the $j^\dagger$-th row of \smash{$M_\lambda^{(i)}$} under $D_i$ coincides with the $j^\dagger$-th row of~\smash{$M_\kappa^{(i)}$}.
\end{proof}





\begin{proof}[Proof of Theorem~\ref{thm:determinant_formula_for_ksmall_g}]
Let $w$ be the element of \smash{$\hat{W}_{G}^0$} associated with $\lambda$.
Since $\lambda$ is $k$-small, there exists a~reduced expression $w=vs_ls_{l+1}\dots s_{n-1}s_0$ for some \label{$w'$ is not affine Grassmannian.}
$v\in \hat{W}_{G}$, where $l$ is the $n$-residue at the bottom of the first column of $\lambda$.


We show the theorem by induction on $\ell(v)$.
The case $\ell(v)=0$ is done, so suppose $\ell(v)>0$.
In this case, $\lambda$ contains at least one removable box
not included in the first column.
Let $i$ be the $n$-residue of the box.
Let $\mu$ the partition obtained from $\la$ by removing
the box of $n$-residue $i$.

(i) If $i\neq 0$, by induction hypothesis, we have
\[
\tilde{g}^{(k)}_\lambda(y|b)=
D_i\bigl(
\tilde{g}^{(k)}_{\mu}(y|b)
\bigr)
=
D_i
\left(
\frac{
\det(M_{\mu})
}{\co_\la(y)}
\right).
\]
Since $\lambda$ is $k$-small, we easily show that $\co_\la(y)$ is invariant under $s_i$.
Therefore, from Lemma~\ref{lem:comparing_determinants}, we have
\[
D_{i}\left(
\frac{\det M_{\mu}}{\co_\la(y)}
\right)
=
\frac{
D_{i}(\det M_{\mu})}{\co_\la(y)}
=
\frac{\det M_{\lambda}}{\co_\la(y)}.
\]

(ii) For $i=0$, because $\Omega(b_1|y)\co_\la(y)$ is $s_\theta$-invariant, we have
\[
\begin{aligned}
\tilde{g}^{(k)}_\lambda(y|b)
&
=
\Omega(b_1|y)
D_\theta
\left(
\frac{
\det M_{\mu}
}{\Omega(b_1|y)\co_\la(y)}
\right)\\
&
=
\Omega(b_1|y)
\frac{
D_\theta(\det M_{\mu})}{\Omega(b_1|y)\co_\la(y)}
=
\frac{\det M_{\lambda}}{\co_\la(y)}
\end{aligned}
\]
from Lemma~\ref{lem:comparing_determinants}.
\end{proof}

\begin{Corollary}\label{cor:g_rectangle}
Suppose $i+j\le n$.
Then we have
\[
\tilde{g}^{(k)}_{(i^j)}(y|b)=
\frac{
\det(M'_{(i^j)})
}{\co_{(i^j)}(y)},\qquad
M'_{(i^j)}=\omega^{-j}
\left(
\begin{matrix}
Z_{[1,j]}^{[1,i+j]}\\\hline
{
\bigl(P^{-1}\bigr)}^{[1,i+j]}_{[j+1,i+j]}
\end{matrix}
\right).
\]
\end{Corollary}


\begin{proof}
Note also that the size of the square matrix $M_{(i^j)}$ is $i+j$.
Let $W$ be the matrix consisting of the last $i$ rows of $M_{(i^j)}$.
Then we can write \begin{equation}
M_{(i^j)}=\left(
\begin{matrix}
\omega^{-j}Z_{[1,j]}^{[1,i+j]}\\\hline
W
\end{matrix}
\right).\label{eq:M rec}
\end{equation}


There is a~lower unitriangular $i\times i$ matrix $N$
such that
\begin{equation}
N \omega^j(W)
=\bigl(P^{-1}\bigr)_{[j+1,i+j]}^{[1,i+j]}.
\label{eq:N omega j W}
\end{equation}
In fact, if we define
$
N:
=
N_1
N_2
\cdots
N_{i-1}$,
where
\smash{$N_m=1-\sum_{j=1}^m{\rm e}^{-a_j}E_{r+j-m,r+j-m-1}$}, then by using the explicit formula for $P^{-1}$, it is straightforward to check~\eqref{eq:N omega j W}.




Combined with~\eqref{eq:M rec} and~\eqref{eq:N omega j W}, we have
\[
\begin{pmatrix}
E_j & O\\
O & N
\end{pmatrix}\omega^j M_{(i^j)}=\left(
\begin{matrix}
Z_{[1,j]}^{[1,i+j]}\\\hline
N\omega^j W
\end{matrix}
\right)=\left(
\begin{matrix}
Z_{[1,j]}^{[1,i+j]}\\\hline
\bigl(P^{-1}\bigr)_{[j+1,i+j]}^{[1,i+j]}
\end{matrix}
\right)=\omega^j M_{(i^j)}'.
\]
Therefore, by taking determinants we
obtain the desired result.
\end{proof}


\begin{Corollary}\label{cor:k rect and tau}
We have
 \begin{align}
 \tilde{g}^{(k)}_{R_i}(y|b)&=
 \frac{\sigma_{n-i}}{\co_{R_i}(y)},\label{eq:det sigma i}
 \\ g^{(k)}_{R_i}(y|b)&=
{\rm e}^{\sum_{s=0}^{n-i-1}a_{i-s}}
 \frac{
 \tau_{n-i}}{\co_{R_i}(y)}.\label{eq:det tau i}\end{align}
\end{Corollary}
\begin{proof}
Let $M'_{R_i}$ be the $n\times n$ matrix in Corollary~\ref{cor:g_rectangle}. We have
\[
\omega^{i}M'_{R_i}=
\left(
\begin{matrix}
Z_{[1,n-i]}^{[1,n]}\\\hline
{
\bigl(P^{-1}\bigr)}^{[1,n]}_{[n-i+1,n]}
\end{matrix}
\right).
\]
So we have
\[
\bigl(\omega^{i}M'_{R_i}\bigr)P=
\left(
\begin{matrix}
(ZP)_{[1,n-i]}^{[1,n]}\\
\hline
O_{n-i}|E_{i}
\end{matrix}
\right).
\]
By taking determinants, we have
\[
\omega^{i}\det(M'_{R_i})
=\det\bigl(\omega^{i}M'_{R_i}\bigr)=\det\left((ZP)^{[1,n-i]}_{[1,n-i]}\right)=\sigma_{n-i}.
\]
Since $\sigma_{n-i}$ is $S_n$-invariant (see Corollary~\ref{cor:tausigma invariant}), we obtain
\eqref{eq:det sigma i}.
%
It is straightforward to obtain~\eqref{eq:det tau i} from~\eqref{eq:det sigma i}
by using~\eqref{eq:sigma(Omega)} and
Proposition~\ref{prop:v1_of_kKShur}.
\end{proof}



\subsection[Determinantal formula for g\^{}(k)\_lambda]{Determinantal formula for $\boldsymbol{g^{(k)}_\lambda}$}

The determinantal formula for the $K$-$k$-Schur function $g^{(k)}_\lambda(y|b)$ is also obtained for a~$k$-small $\lambda$.
\begin{Theorem}\label{thm:determinantal_formula_for_g}
Let $\lambda$ be a~$k$-small partition.
Let $l$, $r$ be as in Theorem~{\rm\ref{thm:determinant_formula_for_ksmall_g}}.
Let $W_\la$ be the matrix consisting of the last $i$ rows of $M_{\la}$.
Then \begin{equation}
g^{(k)}_{\lambda}(y|b)
={\rm e}^{\sum_{x\in \mathrm{diag}(\lambda)}(a_{r(x)+1}-a_{b(x)})}
\frac{\det(N_\lambda)}{\co_\la(y)},\qquad
N_\lambda=
\left(
\begin{matrix}
 (ZA)_{[l,n]}^{[l,n+r]} \\ \hline
 W_{\la}
\end{matrix}
\right).
\label{eq:det g small la}
\end{equation}
\end{Theorem}
\begin{proof} Since $\sigma^{-1}(Z)=ZA$ and $\sigma$ is $\Rep$-linear, we have
$\sigma^{-1}(M_\la)=N_\la.
$ In view of this,
\eqref{eq:det g small la}
is obtained from~\eqref{eq:matrix_M} by applying~\eqref{eq:sigma(Omega)} and Proposition~\ref{prop:v1_of_kKShur}.
\end{proof}


\begin{Example}\label{eq:n6_7boxes} For $\la$ in Example
 \ref{ex:k small det}, $g^{(5)}_{\lambda}(y|b)$ is
\[
\begin{aligned}
&
\dfrac{{\rm e}^{a_2+a_3+a_5}}
{\Omega(b_4|y)\Omega(b_5|y)\Omega(b_6|y)}
\begin{vmatrix}
 z'_{44}&z'_{45}&z'_{46}&z'_{47}&z'_{48}&z'_{49}\\
 0&z'_{55}&z'_{56}&z'_{57}&z'_{58}&z'_{59}\\
 0&0&z'_{66}&z'_{67}&z'_{68}&z'_{69}\\
 -f_3^{\langle 3,1\rangle;\la} & f_2^{\langle 2,1\rangle;\la} & -f_1^{\langle 1,1 \rangle;\la} & 1 & 0 & 0\\
 0&0&f_2^{\langle 2,2\rangle;\la}&-f_1^{\langle 1,2 \rangle;\la} & 1 &0
\end{vmatrix},
\end{aligned}
\]
where
$z'_{ij}={\rm e}^{-a_i}z_{ij}-z_{i,j-1}$.
Note that
$z_{ij}'
$ is specialized to
$h_{j-i}(y)-h_{j-i-1}(y)$
at $a_i=0$.
\end{Example}




\section[k-rectangle factorization property]{$\boldsymbol{k}$-rectangle factorization property}\label{sec:krect}
Define for $1\le j\le n$
\begin{equation}
\label{eq:def de}
\de_j:=
\Omega(b_{n-j+1}|y)^{-1}\cdots \Omega(b_{n}|y)^{-1}.
\end{equation}
Note that for $1\le i\le n-1$, we have
\begin{equation}
\de_{n-i}
=\co_{R_{i}}(y)^{-1}.
\label{eq:d Omega R}
\end{equation}
\begin{Lemma}\label{lem:commutation_relation_Omega}

Let $1\le j\le n$. We consider $\de_j$ a~linear transformation of
$\hat{\Lambda}^{R(T)}$ by multiplication.
For any $i\in \tilde{I}$, we have
\begin{equation}
{D}_i \circ \de_{j}
=\de_j\circ \bigl(\omega^{-j}D_{i+j} \omega^{j}\bigr),
\label{eq:D d}
\end{equation}
where the subscripts are taken modulo $n$.
\end{Lemma}
\begin{proof}
Let us consider the case $i\ne 0$ and $i+j= 0\mod n $.
We have
\begin{align*}
D_i \circ \de_j&=
D_i \circ \Omega(b_{i+1}|y)^{-1}\dots \Omega(b_{n}|y)^{-1}\qquad
\text{by assumption $i+j\equiv 0 \mod n$}\\
&=D_i \circ
\Omega(b_{i}|y)
\Omega(b_{i}|y)^{-1}
\Omega(b_{i+1}|y)^{-1}\de_{j-1}\\
&=\Omega(b_{i}|y)^{-1}
\Omega(b_{i+1}|y)^{-1}\de_{j-1}
\cdot
D_i \circ
\Omega(b_{i}|y)=\de_{j}
\circ
\Omega(b_{i}|y)^{-1}
\circ
D_i \circ
\Omega(b_{i}|y).
\end{align*}
In the third equality, we used that
$\Omega(b_{i}|y)^{-1}
\Omega(b_{i+1}|y)^{-1}\de_{j-1}$
is $s_i$-invariant.
It is straightforward to show
$
\Omega(b_{i}|y)^{-1}
\circ
D_i \circ
\Omega(b_{i}|y)=\omega^i
\circ D_0 \circ \omega^{-i}$.
Hence we have~\eqref{eq:D d} in this case.
The other cases are left to the reader.
\end{proof}

\begin{Theorem}\label{thm:$k$-rect}
Let $\lambda$ be an arbitrary $k$-bounded partition.
Then we have
\[
\tilde{g}^{(k)}_{\lambda\cup R_i}(y|b)=
\tilde{g}^{(k)}_{\lambda}\bigl(y| \omega^{i}b\bigr)\cdot\tilde{g}^{(k)}_{R_i}(y|b).
\]
\end{Theorem}
\begin{proof}
Let $w_\lambda=s_{j_1}s_{j_2}\cdots s_{j_l}$ be a~reduced expression.
Let $w_\la'=s_{j_1+i}s_{j_2+i}\cdots s_{j_l+i}$.
It is shown in \cite[Lemma~2.15]{Tak} that we have a~length additive decomposition
$
w_{\la \cup R_i}=w'_\la w_{R_i}$.

Therefore, we have
\begin{align*}
\tilde{g}^{(k)}_{\lambda\cup R_i}(y|b)&=D_{w'_\la}D_{R_i}(1)=D_{w'_\la}\bigl(\tilde{g}^{(k)}_{R_i}(y|b)\bigr)
\\
&=D_{w'_\la}\left(\frac{\sigma_{n-i}}{\co_{R_i}(y)}\right)\qquad\text{by~\eqref{eq:det sigma i}}\\
&=\sigma_{n-i}\cdot D_{w'_\la}{\co_{R_i}(y)}^{-1}\qquad\text{by Proposition~\ref{prop:S_n_invariant_functions}}
\\
&=\sigma_{n-i}\cdot D_{w'_\la}\de_{n-i}\qquad\text{by~\eqref{eq:d Omega R}}\\
&=\sigma_{n-i} \de_{n-i}\cdot \omega^{i} D_{w_\la}\omega^{-i}(1)
\qquad\text{by Lemma~\ref{lem:commutation_relation_Omega}}
\\
&=\tilde{g}^{(k)}_{R_i}
(u|b)
\cdot \omega^{i}( D_{w_\la}(1))=\tilde{g}^{(k)}_{R_i}
(u|b)
\cdot \tilde{g}^{(k)}_{\la}
\bigl(u| \omega^i b\bigr).\tag*{\qed}
 \end{align*}\renewcommand{\qed}{}
\end{proof}



\section{Proof of Theorem~\ref{thm:max factor}}\label{sec:max factor}



This is obtained by computing the image of $\mathfrak{G}_{w_\circ}^Q$
under $\tilde{\Phi}_n$ in two different ways.

\subsection[Phi\_n(fac\_i)]{$\boldsymbol{\tilde{\Phi}_n(\fac_i)}$}

Recall that $\fac_i$ $(1\le i\le n-1)$ is defined by~\eqref{eq:def phi}.
\begin{Proposition}\label{prop:phi}
For $1\le i\le n-1$,
\[
\tilde{\Phi}_n(\fac_i)=
\frac{\prod_{l=1}^i\Omega(b_{i+l+1}|y)}{\sigma_i}\cdot\tilde{g}_{(n-i-1)^i}^{(k)}\bigl(y|\omega^{2i+1}b\bigr).
\]
\end{Proposition}
Note that for $i=n-1$,
the partition appearing on the right-hand side is $\varnothing$.
Since $\mathfrak{G}_{w_\circ}^Q$ is the product of $\fac_i$ $(1\le i\le n-1)$, we immediately obtain the following formula.
\begin{Corollary}\label{cor:image_of_Groth}
We have
\begin{equation}\label{eq:image_of_Groth}
\tilde{\Phi}_n\bigl(\mathfrak{G}_{w_\circ}^Q\bigr)
=
\frac{\prod_{i=1}^{n-1}\prod_{l=1}^i\Omega(b_{i+l+1}|y)}{{\sigma_1\sigma_2\cdots \sigma_{n-1}}}\cdot
{
\prod_{i=1}^{n-2}
\tilde{g}^{(k)}_{(n-i-1)^i}\bigl(y|\omega^{2i+1}b\bigr)
}.
\end{equation}
\end{Corollary}
Here we note
that the factor
$\prod_{i=1}^{n-1}\prod_{l=1}^i\Omega(b_{i+l+1}|y)$ in~\eqref{eq:image_of_Groth}
can be expressed as follows:
\begin{equation}
\label{eq:Omega_n}
\prod_{i=1}^n\Omega(b_i|y)^m\qquad \text{for} \ n=2m+1,\qquad
\prod_{i=1}^m\Omega(b_{2i-1}|y)^{m}\Omega(b_{2i}|y)^{m-1}\qquad \text{for} \  n=2m.
\end{equation}

As an application of Corollary~\ref{cor:image_of_Groth}, we have a~proof of Proposition~\ref{prop:long_iota_invariant}.
\begin{proof}[Proof of Proposition~\ref{prop:long_iota_invariant}]
The product \smash{$\prod_{i=1}^{n-2}\tilde{g}^{(k)}_{(n-i-1)^i}\bigl(y|\omega^{2i+1}b\bigr)$} is $\iota$-invariant because
\[
\iota\bigl(\tilde{g}^{(k)}_{(n-i-1)^i}\bigl(y|\omega^{2i+1}b\bigr)\bigr)
=
\tilde{g}^{(k)}_{i^{(n-i-1)}}\bigl(y|\omega^{-2i-1}b\bigr)
=\tilde{g}^{(k)}_{i^{(n-i-1)}}\bigl(y|\omega^{2(n-i-1)+1}b\bigr)
\]
 by
Proposition~\ref{prop:iota $K$-k}.
Let $\Omega_n$ denote the factor~\eqref{eq:Omega_n}.
We will show the factor $\Omega_n\cdot (\sigma_1\sigma_2\cdots \sigma_{n-1})^{-1}$ is also $\iota$-invariant.
In fact, from~\eqref{eq:Omega_n} and~\eqref{eq:iota Omega} we have
$
\iota(\Omega_n)
=\Omega_n/ \sigma_n^{n-1}.
$
From~\eqref{prop:S_on_tau},
we have~${
\iota(\sigma_1\cdots \sigma_{n-1})=\sigma_1\cdots \sigma_{n-1}\cdot \sigma_n^{n-1}}$.

Hence, from Corollary~\ref{cor:image_of_Groth},
$\tilde{\Phi}_n\bigl(\mathfrak{G}_{w_\circ}^Q\bigr)$ is invariant under the action of $\iota$.
\end{proof}

The rest of this subsection is devoted to the proof of Proposition~\ref{prop:phi}.


We start with the expression
\[
\tilde{\Phi}_n(\fac_i)=
\frac{\det\left((E-{\rm e}^{a_{i+1}}C_A )P^{-1}ZP\right)_{[1,i]}^{[1,i]}}{\sigma_i}.
\]
This is obtained from
\eqref{eq:phi det} and
\eqref{eq:PhiL_Z}.



Let us study the matrix
$
\Xi:=
(E-{\rm e}^{a_1}C_A)P^{-1}ZP$.
Because the entries of $C_A$ and $P^{-1}ZP$ are $S_n$-invariant (see Remark~\ref{rem:PinvZP}),
we have
$
 \omega^i \Xi=
(E-{\rm e}^{a_{i+1}}C_A)P^{-1}ZP$.

In order to finish the proof of Proposition~\ref{prop:phi}, it suffices to prove the following.
\begin{Lemma}\label{lem:det Xi}
Let \smash{$M_{(n-i-1)^i}'$} be as defined in Corollary~{\rm\ref{cor:g_rectangle}}.
We have
\begin{equation}
\det \Xi_{[1,i]}^{[1,i]}=\omega^{i+1} \det M_{(n-i-1)^i}'.
\label{eq:det Xi}
\end{equation}
\end{Lemma}
\begin{proof}
Recall $C_A=P^{-1}AP$.
Because $P^{-1}$ is lower unitriangular, we have
\[
\det\Xi^{[1,i]}_{[1,i]}
=
\det((E-{\rm e}^{a_1}A)ZP)_{[1,i]}^{[1,i]}.
\]
As all elements in the $1$st column of $E-{\rm e}^{a_1}A$ are zero and \smash{$(E-{\rm e}^{a_1}A)_{[1,n-1]}^{[2,n]}$} is lower triangular, the determinant \smash{$\det((E-{\rm e}^{a_1}A)ZP)_{[1,i]}^{[1,i]}$} decomposes as
\[
\det(E-{\rm e}^{a_1}A)^{[2,i+1]}_{[1,i]}\det(ZP)_{[2,i+1]}^{[1,i]}
={\rm e}^{ia_1}\det(ZP)_{[2,i+1]}^{[1,i]}.
\]
Therefore, we obtain
\begin{align}
\det\Xi^{[1,i]}_{[1,i]}
&={\rm e}^{ia_1}\det(ZP)_{[2,i+1]}^{[1,i]}
={\rm e}^{ia_1}\det\left(
\begin{matrix}
Z_{[2,i+1]}^{[1,n]} \\\hline
\bigl(P^{-1}\bigr)_{[i+1,n]}^{[1,n]}
\end{matrix}
\right), \label{eq:Ni det}
\end{align}
where for the last equality we used $\det(P)=1$ (cf.\ proof of Corollary~\ref{cor:k rect and tau}).
By row reduction, we can replace \smash{$\bigl(P^{-1}\bigr)_{[i+1,n]}^{[1,n]}$} in~\eqref{eq:Ni det} by
\begin{equation}
\left(
\begin{array}{@{}c|c@{}}
 (-{\rm e}^{-a_1})^{i} &\ast\ \cdots\ \ast \\\hline
 \bm{0}_{n-i-1} & \omega\bigl(P^{-1}\bigr)_{[i+1,n-1]}^{[1,n-1]}
\end{array}
\right).
\label{eq:sweep}
\end{equation}
In fact we sweep out the first column of \smash{$\bigl(P^{-1}\bigr)_{[i+1,n]}^{[1,n]}$} taking $(1,1)$ entry $(-{\rm e}^{-a_1})^i$ as
the pivot. By using the following obvious identity
\[
h_m(x_1,\dots,x_n)-h_{m-1}(x_1,\dots,x_n)x_n=h_{m}(x_1,\dots,x_{n-1})
\] repeatedly, we obtain~\eqref{eq:sweep}.
Thus by expanding along the first column, we have
\[
\det\Xi^{[1,i]}_{[1,i]}=
\det\left(
\begin{matrix}
Z_{[2,i+1]}^{[2,n]} \\\hline
\omega\bigl(P^{-1}\bigr)_{[i+1,n-1]}^{[1,n-1]}
\end{matrix}
\right).
\]
If we compare this with
\[M_{(n-i-1)^i}'=
\omega^{-i} \left(
\begin{matrix}
Z^{[1,n-1]}_{[1,i]}\\
\hline
\bigl(P^{-1}\bigr)^{[1,n-1]}_{[i+1,n-1]}
\end{matrix}
\right),
\]
and use~\eqref{eq:pi_zij},
we obtain~\eqref{eq:det Xi}.
\end{proof}

\begin{proof}
[Proof of Proposition~\ref{prop:phi}]
\begin{align*}
\tilde{\Phi}_n(\fac_i)&=
{\sigma_i}^{-1}{\det}\bigl(\omega^{i} \Xi\bigr)^{[1,i]}_{[1,i]}={\sigma_i}^{-1}{\det}\bigl(\omega^{2i+1} M'_{(n-i-1)^i}\bigr)^{[1,i]}_{[1,i]} \qquad\text{by Lemma~\ref{lem:det Xi}}\\
&={\sigma_i}^{-1}
\omega^{2i+1}\bigl({\co_{(n-i-1)^i}(y)}\cdot{\tilde{g}^{(k)}_{(n-i-1)^i}(y|b)}\bigr)\qquad\text{by Corollary~\ref{cor:g_rectangle}}\\
&={\sigma_i}^{-1}\prod_{l=1}^i\Omega(b_{i+l+1}|y)\cdot\tilde{g}^{(k)}_{(n-i-1)^i}\bigl(y|\omega^{2i+1}b\bigr).\tag*{\qed}
\end{align*}\renewcommand{\qed}{}
\end{proof}

\subsection{Proof of Theorem~\ref{thm:max factor}}
\begin{Lemma}\label{lem:rho check}
Let \smash{$\rho^\lor=\sum_{i=1}^{n-1}\varpi_i^\lor$}.
Suppose $n$ is even and $n=2m$. Then
$\rho^\lor+\varpi_m^\lor\in Q^{\lor}$ and
$
w_\circ t_{-\rho^\lor-\varpi_m^\lor}
=x_{R_m\cup \nu_n}$.
Suppose $n$ is odd. Then $\rho^\lor\in Q^{\lor}$, and
$
w_\circ t_{-\rho^\lor}
=x_{\nu_n}$.
\end{Lemma}
\begin{proof}
It is well known that
\begin{equation}\label{eq:rho check}
\rho^\lor=\frac{1}{2}\sum_{i=1}^{n-1
}i(n-i)\alpha_i^\lor.
\end{equation}
We see that if $n$ is odd then $\rho^\lor\in Q^\lor$
and if $n$ is even then
$\rho^\lor+\varpi_m^\lor\in Q^\lor$.


We work in the extended affine Weyl group $\tilde{W}_G:=\langle \pi\rangle\ltimes \hat{W}_G$
where $\pi^n=\mathrm{id}$ and $\pi s_i=s_{i+1}\pi$ for $i\in \tilde{I}=\Z/n\Z$.
By using \cite[Lemma 4.6]{IIN}, we have
\begin{equation}
\label{eq:w0t}
w_\circ t_{-\rho^\lor}=
w_\circ t_{-\sum_{i=1}^{n-1}\varpi_i^\lor}=
\pi^{-\sum_{i=1}^{n-1}i}x_{\nu_n}.
\end{equation}
Suppose $n=2m$ is even.
Then \smash{$
\pi^{-\sum_{i=1}^{n-1}i}=\pi^{m}$}.
From this, together with
\smash{$t_{-\varpi_m^\lor}=\pi^{-m} x_{R_m}$},
we have
\[
w_\circ t_{-\rho^\lor-\varpi_m^\lor}
=\pi^{m}x_{\nu_n}\pi^{-m} x_{R_m}
=x_{\nu_n\cup R_m},
\]
where, for the last equality, we use \cite[Lem\-ma~2.15]{Tak}.
If $n$ is odd, then
\smash{$
\pi^{-\sum_{i=1}^{n-1}i}=\mathrm{id}$} and the desired equality is nothing but~\eqref{eq:w0t}.
\end{proof}


\begin{proof}[Proof of Theorem~\ref{thm:max factor}]
We first note that $\nu_n$
and $R_m$ are invariant under $\omega_k$.

Suppose $n$ is odd.
From Theorem~\ref{thm:main},
\[
\tilde{\Phi}_n\bigl(Q^{-\rho^\lor}\mathfrak{G}_{w_\circ}^Q\bigr)=\tilde{g}_{\nu_n^{\omega_k}}^{(k)}(y|b)=
\tilde{g}_{\nu_n}^{(k)}(y|b).
\]
From
\eqref{eq:rho check}, it is straightforward to show
\[
\tilde{\Phi}_n\bigl(Q^{\rho^\lor}\bigr)=
\frac{\sigma_{n}^{(n-1)/2}}{\sigma_1\sigma_2\cdots \sigma_{n-1}}.
\]
Comparing these equations with
Proposition~\ref{prop:phi}
and using
$\sigma_n=\Omega(b_1|y)\cdots\Omega(b_n|y)$, we obtain~\eqref{eq:odd}.

Next we consider the case
$n=2m$. From Theorem~\ref{thm:main}, we have
\begin{align*}
\tilde{\Phi}_n
\bigl(Q^{-\rho^\lor-\varpi_m^\lor}
\mathfrak{G}_{w_\circ}^Q\bigr)
&=\tilde{g}^{(k)}_{(\nu_n\cup R_m)^{\omega_k}}(y|b)\qquad \text{by Lemma~\ref{lem:rho check}}\\
&=\tilde{g}_{\nu_n\cup R_m}^{(k)}(y|b)=\tilde{g}^{(k)}_{R_m}(y|b)\tilde{g}^{(k)}_{\nu_n}(y|\omega^m b)\qquad\text{by Theorem~\ref{thm:$k$-rect}.}
\end{align*}
From~\eqref{eq:rho check}, it is straightforward to show
\[
\tilde{\Phi}_n\bigl(Q^{\rho^\lor+\varpi_m^\lor}\bigr)=
\frac{\sigma_{n}^m}{\sigma_m(\sigma_1\sigma_2\cdots \sigma_{n-1})}.
\]
Comparing these equations with
Proposition~\ref{prop:phi}, and
using $\sigma_n=\Omega(b_1|y)\cdots\Omega(b_n|y)$ together with
\smash{$\tilde{g}_{R_m}^{(k)}(y|b)=\sigma_m/\co_{R_m}(y)$}, we obtain~\eqref{eq:even}.
\end{proof}

\appendix

\section{Discrete relativistic Toda lattice}\label{sec:Hirotas_bilinear_form}

Let $L=MN^{-1}$ be the Lax matrix of the relativistic Toda equation~\eqref{eq:rel_Toda}.
Define the new matrix $L^+=N^{-1}M$ by switching the position of $M$ and $N^{-1}$.
By the Gauss decomposition~${L^+=M^+\bigl(N^+\bigr)^{-1}}$, we can define the rational map $(N,M)\mapsto \bigl(N^+,M^+\bigr)$, where $N^+$ and~$M^+$ are matrices of the form
\[
N^+=
\begin{pmatrix}
1& \\
-Q^+_1z^+_1 & 1\\
 & \!\!-Q^+_2z^+_2 & \ddots\\
 & & \ddots & 1\\
 & & & -Q^+_{n-1}z^+_{n-1} &1
\end{pmatrix},\qquad
M^+=
\begin{pmatrix}
z^+_1& -1\\
 & z^+_2 &\ddots\\
 & & \ddots & -1\\
 & & & z^+_n
\end{pmatrix}.
\]

Indeed, this correspondence defines an automorphism over the quotient field of~$R[T][z,Q]$, written as
$%\label{eq:dToda}
\mathrm{dToda}\colon \CC(z_i,Q_i)\to \CC(z_i,Q_i)$, $ z_i\mapsto z_i^+$, $ Q_i\mapsto Q_i^+$.
This birational map is explicitly written as
\[
Q^+_i=\frac{z_{i}}{z_{i+1}}Q_i, \qquad
z^+_i=\frac{1-Q^+_{i-1}}{1-Q^+_{i}}z_i
, \qquad Q_0:=0,
\]
which is known as \textit{the discrete Relativistic Toda lattice} \cite{Suris}.


As $NL^+=LN$, the two matrices $L$ and $L^+$ are similar.
Following the same argument as~\eqref{eq:Gauss_decomp} and~\eqref{eq:AZ_to_L}, we have the pair of matrices $U^+,R^+$ satisfying
\[
P^{-1}Z^+AP=\bigl(U^+\bigr)^{-1}R^+\qquad \mbox{and} \qquad L^+=U^+C_A\bigl(U^+\bigr)^{-1}=R^+C_A(R^+)^{-1}.
\]
Comparing them with $L^+=N^{-1}LN=M^{-1}LM$, we obtain the equations
$
U^{+}=N^{-1}U$, $ R^+=c \cdot M^{-1}R
$
with some nonzero constant $c$.
From these expressions, we can show $P^{-1}Z^+AP=c\cdot P^{-1}ZP$, and then,
$%\label{eq:Z+_Z}
Z^+=cZA^{-1}$.
The correspondence $Z\mapsto Z^+$ defines the same map as $\sigma$ (see~\eqref{eq:def_of_sigma}) on the centralizer space $\tilde{\mathscr{Z}}$.
Symbolically, we have $\sigma=\mathrm{dToda}$.

Let
$\tau^+_i:=\sigma_i=\sigma(\tau_i)$
and $\tau^-_i:=\sigma^{-1}(\tau_i)$ be the tau-functions at the next and previous time step, respectively.
Substituting~\eqref{eq:Q_to_tau} and~\eqref{eq:z_to_tau} to
\[
z_i^+=\frac{1-Q^+_{i-1}}{1-Q^+_{i}}z_i,
\] we obtain
\[
\frac{\tau_i^2-\tau_{i+1}\tau_{i-1}}{\tau_i^+\tau_i^-}=
\frac{\tau_{i-1}^2-\tau_{i}\tau_{i-2}}{\tau_{i-1}^+\tau_{i-1}^-}=\cdots
=
\frac{\tau_{1}^2-\tau_{2}\tau_{0}}{\tau_{1}^+\tau_{1}^-}=1.
\]
Finally, we obtain
$%\label{eq:bilinear_form}
\tau_i^2-\tau_{i+1}\tau_{i-1}=\tau_i^+\tau_i^-$,
which is known as \textit{Hirota's bilinear form of the discrete Toda equation}.


\section{Proof of (\ref{eq:polynomialF})}
\label{sec:F}
Note that for any square matrix $A$ of size $n$, we have
\[
\det(\zeta E-A)=
\sum_{i=0}^n(-1)^i\cdot \zeta^{n-i}
\sum_{
\substack{
J\subset [1,n]\\
|J|=i
}
}
\det A_J^J.
\]
Therefore, it suffices to show
\[
\det L_J^J=
\prod_{j\in J,j+1\notin J}(1-Q_j)
\prod_{j\in J}z_j.
\]
Because the matrix $L^{[a,b]}_{[a,b]}$ decomposes as
\[
\begin{pmatrix}
z_a & -1 && & \\
 & z_{a+1}&\ddots& &\\
 & &\ddots& -1 &\\
 & &&\!z_{b-1}&-1\\
 &&&&\!\!\!\!\!(1-Q_b)z_b
\end{pmatrix}
\begin{pmatrix}
1 & & && & \\
-Q_az_a & 1 && & \\
 &\!\!\!\!\!\!\!-Q_{a+1}z_{a+1} &\ddots& &\\
 & & \ddots &1&\\
 &&&\!\!\!\!\!\!\!\!\!-Q_{b-1}z_{b-1}&1
\end{pmatrix}^{-1},
\]
we have
\[
\det L^{[a,b]}_{[a,b]}=
(1-Q_b)z_a\cdots z_b
\qquad
\text{for}\quad 1\le a\le b\le n.
\]
If we decompose
$J$ as
$J=J_1\sqcup \cdots\sqcup J_r$ in such a~way that
$J_i=[a_i,b_i]$
with $a_{i+1}-b_i\ge 2 $, then
$L_J^J$ is blockwise lower triangular with respect to
$J=J_1\sqcup \cdots\sqcup J_r$. So we have
$\det L_J^J
=\prod_{i=1}^r \det L_{J_i}^{J_i}$
which concludes~\eqref{eq:polynomialF}.

\section[Affine K-nil-Hecke action on quantum K-ring]{Affine $\boldsymbol{K}$-nil-Hecke action on quantum $\boldsymbol{ K}$-ring}\label{sec:nil-Hecke_on_QK}
The aim of this appendix is to give a~proof of Proposition~\ref{prop:Quantum_D(1)}. We use the theory of
semi-infinite flag manifolds.

Let $G$ be a~simple, simply-connected complex algebraic group. Let $B$ be a~Borel subgroup of $G$ and $T$ a~maximal torus contained in $B$.
Let $Q$, $P$ be the root lattice and weight lattice, respectively.
Let $W_G=\langle s_i \mid i\in I\rangle$ be the Weyl group of $G$, and $\hat{W}_G=W_G\ltimes Q^\lor$ the affine Weyl group with $\hat{W}_G=\big\langle s_i \mid i\in \Iaf\big\rangle$.
The \emph{semi-infinite flag manifold} associated with $G$
is the (reduced) ind-scheme of ind-infinite type \smash{$\mathbf{Q}_G^{\mathrm{rat}}= G(\mathbb{F})/(T(\C)N(\mathbb{F}))$} (see~\cite{K:Fro,KNS}), where $\mathbb{F}=\C(\!(t)\!)$ and~$N$ is the unipotent radical of $B$. For each $x\in \hat{W}_G$, there is the corresponding Schubert variety~\smash{$\mathbf{Q}_G(x)\subset \mathbf{Q}_G^{\mathrm{rat}}$}. Let $\mathbf{Q}_G=\mathbf{Q}_G(e)$, where $e$ is the identity.
Let $Q^{\lor,+}$ denote the positive part of the coroot lattice.
The $T$-equivariant $K$-group $K_T(\mathbf{Q}_G)$ is, as an $R(T)$-module, the direct product \smash{$\prod_{x\in \hat{W}_{G}^{+}}R(T)\O_{\mathbf{Q}_G(x)}$}, where~${\hat{W}_{G}^{+}=
\big\{wt_\xi\mid w\in W_G,\, \xi\in Q^{\lor,+}\big\}}$.
Kato~\cite{Kato} established an $R(T)$-module isomorphism from
$QK_T(G/B)$ to $K_T(\mathbf{Q}_G)$.
Explicitly,
${\rm e}^\mu \O^wQ^\xi$ in $QK_T(G/B)$, with $\mu\in P$, $w\in W_G$, $\xi \in Q^\lor$,
corresponds to
${\rm e}^{-\mu}\O_{\mathbf{Q}_G(wt_\xi)}$
in $K_T(\mathbf{Q}_G)$; here we follow the convention of \cite[Section~6.1]{LNS:Se} and \cite[Section~5]{MNS1}.


The \emph{semi-infinite length} of $x=wt_\xi\in \hat{W}_G$
is defined by
\smash{$\ell^\sinf(w)
=\ell(w)+2\langle \rho,\xi\rangle$}, where $\rho$ is~the half sum of the positive roots.
Let \smash{$\le_\sinf$} denote the \emph{semi-infinite Bruhat order} (see \cite[Section~2.4]{INS}).
We note that
\[
s_i x>_\sinf x\Longleftrightarrow \ell^\sinf(s_ix)=\ell^\sinf(x)+1, \qquad i\in \tilde{I}
\] (see \cite[Lemma~4.1.2]{INS}).
Let $\hat{W}_G^0$ denote the minimal-length coset representatives for the coset space $\hat{W}_G/W_G$
with respect to the ordinary length function.

The following fact is due to Peterson~\cite{Pet}.
\begin{Lemma}\label{lem:sinf}
For \smash{$x\in \hat{W}_G^0$}, we have
\smash{$\ell(x)=-\ell^{\sinf}(x)$}.
\end{Lemma}
\begin{proof}
Let $x=wt_\lambda\in \hat{W}_G^0$ $\bigl(w\in W_G, \, \lambda\in Q^\vee\bigr)$. Note that $\lambda\in Q^\vee$ is anti-dominant according to \cite[Lemma 3.3]{LS:Acta}.
We have
\begin{align*}
 \ell(x)&=\ell(t_\lambda)-\ell(w)=\langle e\cdot \lambda,-2\rho\rangle-\ell(w)=-\ell^\sinf(x),
\end{align*}
where the first equality uses \cite[Lemma 3.3]{LS:Acta}, and the second equality applies \cite[Lemma 3.2]{LS:Acta} with $w=e$.
\end{proof}


Motivated by the nil-DAHA action on $K_T\bigl(\mathbf{Q}_G^{\mathrm{rat}}\bigr)$
\cite[Theorem 6.5]{KNS} (see also \cite[Section~3.1]{KoNOS}, \cite[Section~2.6]{Orr}),
we define
an endomorphism of
$QK_T(G/B)_Q$,
\[
D_0^Q=T_\theta+
Q^{-\theta^\vee}
\mathscr{O}^{s_\theta}\cdot
s_\theta,
\]
where~$\theta$ is the highest root and $T_\theta$ is defined in~\eqref{eq:Ttheta}.
Then $D_i$ $ \big(i\in \tilde{I}\big)$ satisfy the braid relations, and hence $D_x^Q$ is defined for $x\in \hat{W}_G$.

The next result is transported from
the corresponding one for semi-infinite flag manifolds
(see \cite[Section~2.6]{Orr} and \cite[Section~3.1]{KoNOS}).
\begin{Proposition}\label{prop:left_Dem}
Let $w\in W$. For $i\in I$, we have
\begin{align}
\label{eq:DiQGro}
D_i^Q(\mathscr{O}^w)
&=\begin{cases}
\mathscr{O}^{s_iw} & \text{if}\ s_iw<w,\\
 \mathscr{O}^{w} & \text{if}\ s_iw>w.
\end{cases} \end{align}
Moreover, we have
\begin{align}
\label{eq:D0QGro}
D_0^Q(\mathscr{O}^w)
&=\begin{cases}
 Q^{-w^{-1}(\theta^\vee)}\mathscr{O}^{s_\theta w} & \text{if}\ s_\theta w>w,\\
 \mathscr{O}^{w} & \text{if}\ s_\theta w<w.
\end{cases}
\end{align}
\end{Proposition}
\begin{proof}
Let $x=wt_\xi$
with $w\in W_G$
and $\xi\in Q^{\lor}$.
We use the following fact (see \cite[Appendix~A]{INS})
\begin{align*}
\text{for}\quad i\in I,\qquad s_i x<_\sinf x   \ \Longleftrightarrow \ s_i w< w,\qquad
 s_0 x<_\sinf x \ \Longleftrightarrow \ s_\theta w>w.
\end{align*}
By the isomorphism
$QK_T(G/B)\cong K_T(\mathbf{Q}_G)$ described above,
\eqref{eq:DiQGro} and~\eqref{eq:D0QGro} are obtained from
\cite[equation~(2.28)]{Orr} or \cite[equation~(3.19)]{KoNOS}.
Note that $s_0w=s_\theta w t_{-w^{-1}(\theta^\lor)}$
for $w\in W_G$.
\end{proof}

 \begin{Corollary} \label{cor:BF}
Let $x=wt_\xi\in \hat{W}_G^0$. Then
$
D_x^Q(1)
=Q^\xi \mathscr{O}^{w}$.
\end{Corollary}
\begin{proof}
Let $x=s_{i_1}\cdots s_{i_l}$ be a~reduced expression.
Then we have a~saturated decreasing chain of elements in
$\hat{W}_G^0$
\[
x=s_{i_1}\cdots s_{i_l}
>s_{i_2}\cdots s_{i_l}
>
\cdots
>s_{i_l}>e.
\]
By Lemma~\ref{lem:sinf}, we obtain a~saturated increasing chain of elements
with respect to $<_\sinf$
\[
x=s_{i_1}\cdots s_{i_l}
<_\sinf s_{i_2}\cdots s_{i_l}
<_\sinf
\cdots
<_\sinf s_{i_l}<_\sinf e.
\]
Thus from Proposition~\ref{prop:left_Dem}, we deduce the corollary.
\end{proof}
\begin{proof}[Proof of Proposition~\ref{prop:Quantum_D(1)}]
Apply Corollary~\ref{cor:BF} to the $T$-equivariant
quantum $K$-ring of \linebreak
$\SL_n(\C)$ presented in~\cite{MNS1} and use the fact that
$\Gro{w}(z|\eta)$ represents $\O^w$.
\end{proof}

\begin{thebibliography}{99}
\footnotesize\itemsep=0pt

\bibitem{ACT}
Anderson D., Chen L., Tseng H.H., On the quantum {$K$}-ring of the flag
 manifold, \href{http://arxiv.org/abs/1711.08414}{arXiv:1711.08414}.

\bibitem{BMS}
Blasiak J., Morse J., Seelinger G.H., {$K$}-theoretic {C}atalan functions,
 \href{https://doi.org/10.1016/j.aim.2022.108421}{\textit{Adv. Math.}}
 \textbf{404} (2022), 108421, 39~pages,
 \href{http://arxiv.org/abs/2010.01759}{arXiv:2010.01759}.

\bibitem{CL}
Chow C.H., Leung N.C., Quantum {$K$}-theory of {$G/P$} and {$K$}-homology of
 affine {G}rassmannian,
 \href{http://arxiv.org/abs/2201.12951}{arXiv:2201.12951}.

\bibitem{Gi}
Ginzburg V.A., Perverse sheaves on a~Loop group and Langlands' duality,
 \href{http://arxiv.org/abs/alg-geom/9511007}{arXiv:alg-geom/9511007}.

\bibitem{GK}
Givental A., Kim B., Quantum cohomology of flag manifolds and {T}oda lattices,
 \href{https://doi.org/10.1007/BF02101846}{\textit{Comm. Math. Phys.}}
 \textbf{168} (1995), 609--641,
 \href{http://arxiv.org/abs/hep-th/9312096}{arXiv:hep-th/9312096}.

\bibitem{GL}
Givental A., Lee Y.-P., Quantum {$K$}-theory on flag manifolds,
 finite-difference {T}oda lattices and quantum groups,
 \href{https://doi.org/10.1007/s00222-002-0250-y}{\textit{Invent. Math.}}
 \textbf{151} (2003), 193--219,
 \href{http://arxiv.org/abs/math.AG/0108105}{arXiv:math.AG/0108105}.

\bibitem{HK}
Huq-Kuruvilla I., Relations in twisted quantum {$K$}-rings,
 \href{http://arxiv.org/abs/2406.00916}{arXiv:2406.00916}.

\bibitem{IIM}
Ikeda T., Iwao S., Maeno T., Peterson isomorphism in {$K$}-theory and
 relativistic {T}oda lattice,
 \href{https://doi.org/10.1093/imrn/rny051}{\textit{Int. Math. Res. Not.}}
 \textbf{2020} (2020), 6421--6462,
 \href{http://arxiv.org/abs/1703.08664}{arXiv:1703.08664}.

\bibitem{IIN}
Ikeda T., Iwao S., Naito S., Closed {$k$}-{S}chur {K}atalan functions as
 {$K$}-homology {S}chubert representatives of the affine {G}rassmannian,
 \href{https://doi.org/10.1090/btran/184}{\textit{Trans. Amer. Math. Soc.
 Ser.~B}} \textbf{11} (2024), 667--702,
 \href{http://arxiv.org/abs/2203.14483}{arXiv:2203.14483}.

\bibitem{ISY}
Ikeda T., Shimozono M., Yamaguchi K., Equivariant {$K$}-homology of affine
 {G}rassmannian and {$K$}-theoretic double {$k$}-{S}chur functions,
 \href{https://doi.org/10.1016/j.aim.2026.110894}{\textit{Adv. Math.}}
 \textbf{492} (2026), 110894, 57~pages,
 \href{http://arxiv.org/abs/2026.11089}{arXiv:2026.11089}.

\bibitem{INS}
Ishii M., Naito S., Sagaki D., Semi-infinite {L}akshmibai--{S}eshadri path
 model for level-zero extremal weight modules over quantum affine algebras,
 \href{https://doi.org/10.1016/j.aim.2015.11.037}{\textit{Adv. Math.}}
 \textbf{290} (2016), 967--1009,
 \href{http://arxiv.org/abs/1402.3884}{arXiv:1402.3884}.

\bibitem{K:Fro}
Kato S., Frobenius splitting of {S}chubert varieties of semi-infinite flag
 manifolds, \href{https://doi.org/10.1017/fmp.2021.5}{\textit{Forum Math. Pi}}
 \textbf{9} (2021), e5, 56~pages,
 \href{http://arxiv.org/abs/1810.07106}{arXiv:1810.07106}.

\bibitem{Kato}
Kato S., Loop structure on equivariant {$K$}-theory of semi-infinite flag
 manifolds, \href{https://doi.org/10.4007/annals.2025.202.3.2}{\textit{Ann. of
 Math.}} \textbf{202} (2025), 1001--1075,
 \href{http://arxiv.org/abs/1805.01718}{arXiv:1805.01718}.

\bibitem{KNS}
Kato S., Naito S., Sagaki D., Equivariant {$K$}-theory of semi-infinite flag
 manifolds and the {P}ieri--{C}hevalley formula,
 \href{https://doi.org/10.1215/00127094-2020-0015}{\textit{Duke Math.~J.}}
 \textbf{169} (2020), 2421--2500,
 \href{http://arxiv.org/abs/1702.02408}{arXiv:1702.02408}.

\bibitem{Kim}
Kim B., Quantum cohomology of flag manifolds~{$G/B$} and quantum {T}oda
 lattices, \href{https://doi.org/10.2307/121021}{\textit{Ann. of Math.}}
 \textbf{149} (1999), 129--148,
 \href{http://arxiv.org/abs/alg-geom/9607001}{arXiv:alg-geom/9607001}.

\bibitem{KPSZ}
Koroteev P., Pushkar P.P., Smirnov A.V., Zeitlin A.M., Quantum {$K$}-theory of
 quiver varieties and many-body systems,
 \href{https://doi.org/10.1007/s00029-021-00698-3}{\textit{Selecta
 Math.~(N.S.)}} \textbf{27} (2021), 87, 40~pages,
 \href{http://arxiv.org/abs/1705.10419}{arXiv:1705.10419}.

\bibitem{KoNOS}
Kouno T., Naito S., Orr D., Sagaki D., Inverse {$K$}-{C}hevalley formulas for
 semi-infinite flag manifolds,~{I}: minuscule weights in {ADE} type,
 \href{https://doi.org/10.1017/fms.2021.45}{\textit{Forum Math. Sigma}}
 \textbf{9} (2021), e51, 25~pages,
 \href{http://arxiv.org/abs/2008.10483}{arXiv:2008.10483}.

\bibitem{LLMS}
Lam T., Li C., Mihalcea L.C., Shimozono M., A conjectural {P}eterson
 isomorphism in {$K$}-theory,
 \href{https://doi.org/10.1016/j.jalgebra.2018.07.029}{\textit{J.~Algebra}}
 \textbf{513} (2018), 326--343,
 \href{http://arxiv.org/abs/1705.03435}{arXiv:1705.03435}.

\bibitem{LSS:K}
Lam T., Schilling A., Shimozono M., {$K$}-theory {S}chubert calculus of the
 affine {G}rassmannian,
 \href{https://doi.org/10.1112/S0010437X09004539}{\textit{Compos. Math.}}
 \textbf{146} (2010), 811--852,
 \href{http://arxiv.org/abs/0901.1506}{arXiv:0901.1506}.

\bibitem{LS:Acta}
Lam T., Shimozono M., Quantum cohomology of~{$G/P$} and homology of affine
 {G}rassmannian, \href{https://doi.org/10.1007/s11511-010-0045-8}{\textit{Acta
 Math.}} \textbf{204} (2010), 49--90,
 \href{http://arxiv.org/abs/0705.1386}{arXiv:0705.1386}.

\bibitem{LS:double-Kostant}
Lam T., Shimozono M., From double quantum {S}chubert polynomials to $k$-double
 Schur functions via the {T}oda lattice,
 \href{http://arxiv.org/abs/1109.2193}{arXiv:1109.2193}.

\bibitem{LS10}
Lam T., Shimozono M., From quantum {S}chubert polynomials to {$k$}-{S}chur
 functions via the {T}oda lattice,
 \href{https://doi.org/10.4310/MRL.2012.v19.n1.a7}{\textit{Math. Res. Lett.}}
 \textbf{19} (2012), 81--93,
 \href{http://arxiv.org/abs/1010.4047}{arXiv:1010.4047}.

\bibitem{LM}
Lenart C., Maeno T., Quantum {G}rothendieck polynomials,
 \href{http://arxiv.org/abs/math.CO/0608232}{arXiv:math.CO/0608232}.

\bibitem{LNS:Se}
Lenart C., Naito S., Sagaki D., A general {C}hevalley formula for semi-infinite
 flag manifolds and quantum {$K$}-theory,
 \href{https://doi.org/10.1007/s00029-024-00924-8}{\textit{Selecta
 Math.~(N.S.)}} \textbf{30} (2024), 39, 44~pages,
 \href{http://arxiv.org/abs/2010.06143}{arXiv:2010.06143}.

\bibitem{Mac}
Macdonald I.G., Notes on {S}chubert polynomials, Publications du LaCIM,
 Universit\'e du Qu\'ebec \`a Montr\'eal, 1991, available at \url{https://lacim.uqam.ca/les-parutions/LACIM-Publications-Volume-06.pdf}.

\bibitem{MNS1}
Maeno T., Naito S., Sagaki D., A presentation of the torus-equivariant quantum
 {$K$}-theory ring of flag manifolds of type~{$A$}, {P}art~{I}: {T}he defining
 ideal, \href{https://doi.org/10.1112/jlms.70095}{\textit{J.~Lond. Math.
 Soc.}} \textbf{111} (2025), e70095, 43~pages,
 \href{http://arxiv.org/abs/2302.09485}{arXiv:2302.09485}.

\bibitem{MNS2}
Maeno T., Naito S., Sagaki D., A presentation of the torus-equivariant quantum
 {$K$}-theory ring of flag manifolds of type~{$A$}, {P}art~{II}: quantum
 double {G}rothendieck polynomials,
 \href{https://doi.org/10.1017/fms.2024.147}{\textit{Forum Math. Sigma}}
 \textbf{13} (2025), e19, 26~pages,
 \href{http://arxiv.org/abs/2305.17685}{arXiv:2305.17685}.

\bibitem{Orr}
Orr D., Equivariant {$K$}-theory of the semi-infinite flag manifold as a~nil-{DAHA} module,
 \href{https://doi.org/10.1007/s00029-023-00848-9}{\textit{Selecta
 Math.~(N.S.)}} \textbf{29} (2023), 45, 26~pages,
 \href{http://arxiv.org/abs/2001.03490}{arXiv:2001.03490}.

\bibitem{Pet}
Peterson D., {MIT} lecture notes, 1997.

\bibitem{Rui}
Ruijsenaars S.N.M., Relativistic {T}oda systems,
 \href{https://doi.org/10.1007/BF02097366}{\textit{Comm. Math. Phys.}}
 \textbf{133} (1990), 217--247.

\bibitem{Suris}
Suris Yu.B., A discrete-time relativistic {T}oda lattice,
 \href{https://doi.org/10.1088/0305-4470/29/2/022}{\textit{J.~Phys.~A}}
 \textbf{29} (1996), 451--465,
 \href{http://arxiv.org/abs/solv-int/9510007}{arXiv:solv-int/9510007}.

\bibitem{Tak}
Takigiku M., A {P}ieri formula and a~factorization formula for sums of
 {$K$}-theoretic {$k$}-{S}chur functions,
 \href{https://doi.org/10.5802/alco.45}{\textit{Algebr. Comb.}} \textbf{2}
 (2019), 447--480, \href{http://arxiv.org/abs/1802.06335}{arXiv:1802.06335}.

\end{thebibliography}
\end{document}
